\documentclass[11pt,a4paper]{report}

% ------------------------------------------------------------
% Curso de Aprendizaje por Refuerzo: De la Probabilidad a las Redes Neuronales
% Manual Completo del Estudiante (36 Horas)
% Formato: Libro / Documento PDF de Alta Calidad
% ------------------------------------------------------------

\usepackage[T1]{fontenc}
\usepackage[utf8]{inputenc}
\usepackage[spanish,es-nodecimaldot]{babel}
\usepackage[scaled=0.95]{helvet}
\renewcommand{\familydefault}{\sfdefault}
\usepackage{microtype}
\usepackage{amsmath,amssymb,mathtools}
\usepackage{booktabs,array}
\usepackage{geometry}
\geometry{
  top=2.8cm,
  bottom=2.8cm,
  left=2.7cm,
  right=2.7cm,
  headheight=32pt
}
\usepackage{fancyhdr}
\usepackage{titlesec}
\usepackage{listings}
\usepackage{tikz}
\usetikzlibrary{arrows.meta,calc,positioning,shapes.geometric}
\usepackage[most]{tcolorbox}
\usepackage{hyperref}

% ---------- Paleta de colores unificada ----------
\definecolor{Ink}{HTML}{102A43}       % Texto principal y títulos
\definecolor{Paper}{HTML}{FFFFFF}     % Fondo de página
\definecolor{Mist}{HTML}{D9E2EC}      % Gris azulado para bordes y separadores
\definecolor{Teal}{HTML}{005EA8}      % Azul primario institucional
\definecolor{TealDark}{HTML}{003E73}  % Azul profundo
\definecolor{Coral}{HTML}{D9480F}     % Naranja/coral de énfasis e ideas clave
\definecolor{Gold}{HTML}{B77900}      % Ámbar para recompensas y alertas
\definecolor{Leaf}{HTML}{2F6F44}      % Verde para ejemplos y resultados
\definecolor{CodeBg}{HTML}{0B1D33}    % Fondo oscuro de código
\definecolor{SoftText}{HTML}{486581}  % Texto secundario y notas
\definecolor{LightBg}{HTML}{F0F4F8}   % Fondo tenue para cajas de texto

% ---------- Configuración de hipervínculos ----------
\hypersetup{
  colorlinks=true,
  linkcolor=Teal,
  citecolor=Leaf,
  urlcolor=Coral,
  pdftitle={Aprendizaje por Refuerzo: De la Probabilidad a las Redes Neuronales},
  pdfauthor={Curso Especializado de RL}
}

% ---------- Encabezados y pies de página ----------
\pagestyle{fancy}
\fancyhf{}
\renewcommand{\headrulewidth}{0.6pt}
\renewcommand{\footrulewidth}{0.4pt}
\renewcommand{\headrule}{\hbox to\headwidth{\color{Mist}\leaders\hrule height \headrulewidth\hfill}}
\renewcommand{\footrule}{\hbox to\headwidth{\color{Mist}\leaders\hrule height \footrulewidth\hfill}}

% Formato limpio de marca de capítulo (sin mayúsculas forzadas que saturan el ancho)
\renewcommand{\chaptermark}[1]{\markboth{Capítulo \thechapter.\ #1}{}}

% Encabezados con anchos delimitados para evitar cualquier sobreposición de texto
\fancyhead[L]{\parbox[b]{0.32\linewidth}{\raggedright\footnotesize\color{SoftText}\textit{Aprendizaje por Refuerzo}}}
\fancyhead[R]{\parbox[b]{0.65\linewidth}{\raggedleft\footnotesize\color{Teal}\bfseries\nouppercase{\leftmark}}}

\fancyfoot[C]{\small\color{Ink}\thepage}
\fancyfoot[R]{\tikz[baseline=-0.5ex]{\fill[Teal](0,0) circle (1.5pt);\fill[Coral](4pt,0) circle (1.5pt);\fill[Gold](8pt,0) circle (1.5pt);}}

% ---------- Estilo de Capítulos y Secciones ----------
\titleformat{\chapter}[display]
  {\normalfont\sffamily\bfseries\color{TealDark}}
  {\filleft\fontsize{48}{54}\selectfont\color{Mist}\thechapter}
  {1ex}
  {\titlerule[1.5pt]\vspace{1ex}\filright\fontsize{22}{26}\selectfont\color{Ink}}
  [\vspace{1ex}{\titlerule[0.8pt]}]

\titleformat{\section}
  {\normalfont\Large\bfseries\color{Teal}}
  {\thesection}{1em}{}[\color{Mist}\titlerule]

\titleformat{\subsection}
  {\normalfont\large\bfseries\color{Ink}}
  {\thesubsection}{0.8em}{}

% ---------- Badges de conceptos de RL ----------
\newcommand{\state}[1]{\tikz[baseline=(x.base)]\node[rounded corners=2pt,fill=Teal!13,text=TealDark,inner xsep=5pt,inner ysep=2pt] (x) {\scriptsize\bfseries ESTADO\enspace #1};}
\newcommand{\rlaction}[1]{\tikz[baseline=(x.base)]\node[rounded corners=2pt,fill=Coral!13,text=Coral,inner xsep=5pt,inner ysep=2pt] (x) {\scriptsize\bfseries ACCIÓN\enspace #1};}
\newcommand{\reward}[1]{\tikz[baseline=(x.base)]\node[rounded corners=2pt,fill=Gold!22,text=Ink,inner xsep=5pt,inner ysep=2pt] (x) {\scriptsize\bfseries RECOMPENSA\enspace #1};}

% ---------- Entornos de Cajas de Énfasis (tcolorbox) ----------
\newtcolorbox{ideaclave}[1][]{
  colback=Coral!7,
  colframe=Coral,
  arc=3pt,
  boxrule=1.2pt,
  left=12pt,
  right=12pt,
  top=8pt,
  bottom=8pt,
  title={\textbf{\textcolor{Coral}{Idea Clave}}},
  fonttitle=\sffamily\bfseries,
  coltitle=Coral,
  #1
}

\newtcolorbox{optibot}[1][]{
  colback=Teal!6,
  colframe=Teal,
  arc=3pt,
  boxrule=1.2pt,
  left=12pt,
  right=12pt,
  top=8pt,
  bottom=8pt,
  title={\textbf{\textcolor{Teal}{OptiBot en Acción: Caso Conductor}}},
  fonttitle=\sffamily\bfseries,
  coltitle=Teal,
  #1
}

\newtcolorbox{calculomanual}[1][]{
  colback=Leaf!6,
  colframe=Leaf,
  arc=3pt,
  boxrule=1.2pt,
  left=12pt,
  right=12pt,
  top=8pt,
  bottom=8pt,
  title={\textbf{\textcolor{Leaf}{Cálculo Manual Paso a Paso}}},
  fonttitle=\sffamily\bfseries,
  coltitle=Leaf,
  #1
}

\newtcolorbox{notabox}[1][]{
  colback=Mist!25,
  colframe=SoftText,
  arc=3pt,
  boxrule=0.8pt,
  left=10pt,
  right=10pt,
  top=6pt,
  bottom=6pt,
  title={\textbf{\textcolor{SoftText}{Nota de Diseño y Conexión}}},
  fonttitle=\sffamily\small\bfseries,
  coltitle=SoftText,
  #1
}

% ---------- Estilo de Código Python ----------
\lstdefinestyle{rlcode}{
  language=Python,
  backgroundcolor=\color{CodeBg},
  basicstyle=\ttfamily\footnotesize\color{white!92},
  keywordstyle=\bfseries\color{Gold},
  commentstyle=\itshape\color{Leaf!85},
  stringstyle=\color{Coral!85},
  numbers=left,
  numberstyle=\tiny\color{white!40},
  numbersep=8pt,
  frame=none,
  showstringspaces=false,
  breaklines=true,
  tabsize=2,
  xleftmargin=1.8em,
  aboveskip=10pt,
  belowskip=10pt
}

% ---------- Metadatos del documento ----------
\title{
  \vspace{-1.5cm}
  {\fontsize{13}{16}\selectfont\bfseries\color{Teal}PROGRAMA ACADÉMICO TEÓRICO--PRÁCTICO (36 HORAS)}\\[1.2cm]
  {\fontsize{28}{34}\selectfont\bfseries\color{Ink}Aprendizaje por Refuerzo}\\[0.4cm]
  {\fontsize{16}{20}\selectfont\color{SoftText}De la Probabilidad a las Redes Neuronales Profundas}\\[1cm]
  \tikz{\fill[Teal](0,0) circle (3pt);\fill[Coral](8pt,0) circle (3pt);\fill[Gold](16pt,0) circle (3pt);}
}
\author{\textbf{Profesor Especialista en Aprendizaje por Refuerzo}\\[0.2cm]\small Metodología: Lápiz, Papel y Código $\cdot$ Caso Conductor: OptiBot}
\date{Edición 2026}

\begin{document}

\maketitle

\chapter*{Prólogo y Guía Pedagógica}
\addcontentsline{toc}{chapter}{Prólogo y Guía Pedagógica}

Bienvenido al texto integral del curso \textbf{Aprendizaje por Refuerzo: De la Probabilidad a las Redes Neuronales Profundas}. 

Este curso ha sido diseñado desde cero bajo una filosofía pedagógica estricta: \textbf{cualquier persona con conocimientos elementales de programación en Python y álgebra básica puede y debe comprender la teoría matemática y algorítmica del Aprendizaje por Refuerzo (RL) con total claridad, rigor y profundidad}.

\section*{La Metodología: Lápiz, Papel y Código}
Cada concepto fundamental se aborda en tres fases secuenciales e inseparables:
\begin{enumerate}
  \item \textbf{Intuición conceptual:} ¿Qué problema físico o lógico intentamos resolver?
  \item \textbf{Cálculo manual en la pizarra (Lápiz y Papel):} Desarrollamos los números a mano para un caso reducido, asegurando que las variables y fórmulas no sean una ``caja negra''.
  \item \textbf{Implementación algorítmica (Código Python):} Traducimos la matemática a scripts concisos en NumPy o PyTorch, verificando que la simulación por computadora reproduzca con exactitud el resultado teórico.
\end{enumerate}

\section*{El Hilo Conductor: OptiBot en el Almacén Inteligente}
Para evitar la confusión clásica provocada por cambiar de ejemplo en cada lección, a lo largo de todo el libro seguiremos a un único protagonista: \textbf{OptiBot}, un robot móvil autónomo que opera en un centro de distribución y logística. A medida que avanza el curso, el robot se enfrenta a problemas progresivamente más complejos:
\begin{itemize}
  \item \textbf{Capítulo 0:} Cuantifica la incertidumbre estocástica de los tiempos de entrega.
  \item \textbf{Capítulo 1:} Elige entre 4 empresas de mensajería (Bandidos Multi-Brazo).
  \item \textbf{Capítulo 2:} Modela la degradación aleatoria de la cinta transportadora (Cadenas de Markov).
  \item \textbf{Capítulo 3:} Navega entre zonas del almacén con recompensas y costos (MDP y Ecuaciones de Bellman).
  \item \textbf{Capítulo 4:} Calcula la ruta perfecta conociendo los planos exactos del almacén (Programación Dinámica).
  \item \textbf{Capítulo 5:} Aprende a evaluar rutas promediando turnos completos sin planos (Métodos Monte Carlo).
  \item \textbf{Capítulo 6:} Aprende paso a paso tras cada movimiento individual (SARSA y Q-Learning).
  \item \textbf{Capítulo 7:} Conecta sensores continuos (LiDAR, batería en \%, velocidad) sustituyendo las tablas por Redes Neuronales (Deep Q-Networks).
  \item \textbf{Capítulo 8:} Se adapta a un almacén dinámico con obstáculos imprevistos y proyecta su aprendizaje hacia las fronteras del RL moderno (Policy Gradients, PPO, RLHF).
\end{itemize}

\newpage
\tableofcontents
\newpage

% Inclusión de capítulos auxiliares
\chapter{Probabilidad y Pensamiento Estocástico para RL}
\label{cap:probabilidad}

\section{¿Por qué la programación tradicional no basta?}
En la programación informática convencional, las cosas ocurren de manera predecible y ordenada. Si se escribe una instrucción para sumar dos cantidades o para aplicar un descuento del 15\%, el resultado será siempre exacto y reproducible. Si ese código se ejecuta diez veces o un millón de veces, la computadora devolverá con certeza absoluta la misma respuesta.

Sin embargo, cuando una máquina o un robot intenta tomar decisiones en el mundo real, esa certeza desaparece por completo:
\begin{itemize}
  \item Una rueda puede patinar si el suelo tiene un poco de polvo, humedad o desgaste.
  \item Un sensor de distancia puede reportar dos metros y, un instante después, indicar dos metros y diez centímetros debido a un simple reflejo de luz.
  \item Un trayecto que usualmente toma tres minutos puede demorar diez si alguien dejó una caja bloqueando el pasillo.
\end{itemize}

Si se intentara programar una regla fija para cada uno de estos imprevistos, el código resultante sería infinito, frágil e inmanejable. Es aquí donde el Aprendizaje por Refuerzo recurre a la \textbf{probabilidad}: no como una colección de fórmulas abstractas para memorizar, sino como una brújula práctica que le permite a una máquina convivir con el azar, medir el riesgo y tomar decisiones inteligentes aunque no pueda predecir el futuro inmediato con exactitud.

\begin{optibot}
A lo largo del curso se seguirá a \textbf{OptiBot}, un robot encargado de transportar paquetes desde el área de empaque hasta los muelles de carga en un almacén.

El entorno es cambiante: transitan personas, cruzan otros vehículos y surgen obstáculos imprevistos. Un programa con reglas fijas fallaría ante cualquier imprevisto. Para tomar decisiones acertadas, OptiBot necesita estimar probabilidades: saber qué eventos son comunes, cuáles representan riesgos y qué resultados esperar en promedio.
\end{optibot}

\newpage
\section{Variables Aleatorias: Ponerle números al azar}
En el día a día ocurren eventos cuyo resultado final no se puede predecir con seguridad: el clima de mañana, el tiempo exacto que tardará un envío o la cara que mostrará una moneda al caer.

Para que una computadora pueda trabajar con esa incertidumbre, necesita transformar esos posibles resultados en valores numéricos. A esta herramienta se le denomina \textbf{variable aleatoria discreta} (representada convencionalmente con una letra mayúscula, como $X$):
\begin{quote}
  \textit{Una variable aleatoria es simplemente una regla que asigna un número a cada resultado posible de una situación incierta.}
\end{quote}

Por ejemplo, si se observa el tiempo de viaje que le toma a OptiBot llevar un paquete al muelle de carga, la variable aleatoria $T$ podría tomar los valores discretos de $1$, $2$ o $3$ ($t$) minutos.

A cada uno de esos valores posibles se le asocia su \textbf{probabilidad} de ocurrir, denotada formalmente como $P(X = x)$ o $P(T = t)$ en el caso del tiempo de viaje. Para que esta distribución tenga sentido lógico, debe cumplir con dos condiciones de sentido común:
\begin{enumerate}
  \item \textbf{Las probabilidades se ubican siempre entre 0 y 1} (o entre 0\% y 100\%): No existen probabilidades negativas ni superiores al cien por ciento ($0 \le P(X = x) \le 1$). Si algo es imposible su probabilidad es 0; si es una certeza total, su probabilidad es 1.
  \item \textbf{La suma de todos los casos posibles es igual a 1}: Si se toman en cuenta todos los resultados que pueden llegar a presentarse, es un hecho indiscutible que alguno de ellos terminará ocurriendo:
  \begin{equation}
    \sum_{x} P(X = x) = 1
  \end{equation}
\end{enumerate}

\subsection*{Ejemplo: La distribución de tiempos de viaje de OptiBot}
Para ilustrar estos conceptos de forma concreta, supóngase que se cronometran $100$ viajes de OptiBot transportando paquetes hacia el muelle de carga. Al clasificar los tiempos registrados, se obtiene el siguiente desglose:

\begin{center}
\begin{tabular}{ccc}
  \toprule
  \textbf{Tiempo de viaje ($t$)} & \textbf{Viajes observados} & \textbf{Probabilidad $P(T = t)$} \\
  \midrule
  $t = 1$ minuto  & 20 de 100 viajes & $P(T = 1) = 0.20$ \ ($20\%$) \\
  $t = 2$ minutos & 50 de 100 viajes & $P(T = 2) = 0.50$ \ ($50\%$) \\
  $t = 3$ minutos & 30 de 100 viajes & $P(T = 3) = 0.30$ \ ($30\%$) \\
  \bottomrule
\end{tabular}
\end{center}

A partir de estos datos, se comprueban inmediatamente las dos propiedades obligatorias:
\begin{enumerate}
  \item Cada probabilidad es un valor real entre 0 y 1 ($0.20$, $0.50$ y $0.30$).
  \item La suma de todas las alternativas posibles totaliza la certeza absoluta ($1.00$ o $100\%$):
  \begin{equation*}
    P(T = 1) + P(T = 2) + P(T = 3) = 0.20 + 0.50 + 0.30 = 1.00
  \end{equation*}
\end{enumerate}

Esta tabla describe la \textbf{distribución de probabilidad} de la variable $T$. Aunque no se puede adivinar con exactitud cuántos minutos tomará el próximo viaje en particular, el sistema ahora sabe que lo más habitual es tardar 2 minutos ($50\%$ de probabilidad), mientras que los viajes veloces de 1 minuto ocurrirán solo una de cada cinco veces ($20\%$).

\section{Probabilidad Condicional: Aprender con nueva información}
Las probabilidades rara vez permanecen fijas en el tiempo. Casi siempre se modifican en cuanto se obtiene nueva información del entorno.

Considérese un ejemplo cotidiano del almacén: ¿cuál es la probabilidad de que una rueda de OptiBot resbale? En condiciones normales es muy baja, tal vez de un $2\%$. Pero si los sensores detectan que el piso acaba de ser limpiado con jabón, esa probabilidad se incrementa de inmediato a un $40\%$. La probabilidad cambió radicalmente porque se obtuvo un dato clave sobre la situación actual del entorno.

Esta noción se expresa a través de la \textbf{probabilidad condicional}. La probabilidad de que ocurra un evento $B$, dado que ya ocurrió o se conoce la ocurrencia previa del evento $A$, se escribe $P(B \mid A)$ y se calcula mediante:
\begin{equation}
  P(B \mid A) = \frac{P(A \cap B)}{P(A)}
\end{equation}
Donde $P(A \cap B)$ representa la probabilidad de que ambos eventos ocurran de manera conjunta. En el Aprendizaje por Refuerzo, esta fórmula es la piedra angular de todo el razonamiento: el agente se pregunta en cada paso: \textit{«¿Cuál es la probabilidad de alcanzar una meta exitosa ($B$), dado que se decidió ejecutar la acción de girar a la izquierda ($A$)?»}.

\subsection*{Ejemplo: Cálculo numérico de probabilidad condicional}
Para observar cómo opera esta fórmula en la práctica, considérese un registro histórico de $100$ despachos realizados por OptiBot:
\begin{itemize}
  \item Se define el evento $A$: \textit{«El robot toma la ruta del pasillo central»}. Esto ocurrió en $40$ de los $100$ viajes, por lo que $P(A) = \frac{40}{100} = 0.40$ ($40\%$).
  \item Se define el evento $B$: \textit{«El viaje sufre una demora considerable»}.
  \item Se constató que en $16$ de esos $100$ viajes el robot tomó el pasillo central \textbf{y además} sufrió la demora. Es decir, la probabilidad conjunta es $P(A \cap B) = \frac{16}{100} = 0.16$ ($16\%$).
\end{itemize}

Si OptiBot decide ingresar al pasillo central, ¿cuál es la probabilidad de que termine demorándose? Aplicando la definición de probabilidad condicional:
\begin{equation*}
  P(B \mid A) = \frac{P(A \cap B)}{P(A)} = \frac{0.16}{0.40} = 0.40 \quad (40\%)
\end{equation*}

\textbf{Interpretación práctica:} En todo el almacén, la probabilidad general de sufrir demoras tal vez sea de apenas un $24\%$. No obstante, en cuanto se confirma que el robot eligió el pasillo central (evento $A$), el riesgo de demora se eleva automáticamente al $40\%$. Esta actualización de probabilidades es la que le permite al agente aprender qué acciones son seguras y cuáles aumentan el riesgo de penalización.

\subsection{La Ley de Probabilidad Total}
En muchas ocasiones se desea determinar la probabilidad general de que ocurra un suceso, sabiendo que este puede ocurrir a través de distintas rutas o circunstancias intermedias.

Si OptiBot puede desplazarse hacia el muelle de carga eligiendo entre tres rutas alternativas como el pasillo central, el pasillo norte o el pasillo sur; la probabilidad total de que ocurra un evento (como sufrir una demora) se obtiene evaluando cada alternativa por separado y sumando sus contribuciones ponderadas:
\begin{equation}
  P(B) = \sum_{i=1}^n P(B \mid A_i) \cdot P(A_i)
\end{equation}
En términos prácticos: se calcula la probabilidad de que ocurra el evento en cada camino posible, se multiplica por la frecuencia con la que se transita por dicho camino, y se suman todas esas partes.

\subsection*{Ejemplo: Probabilidad total de sufrir demoras en el almacén}
Siguiendo con el caso de OptiBot, supóngase que el robot dispone de tres rutas posibles hacia el muelle de carga ($A_1$, $A_2$ y $A_3$):
\begin{itemize}
  \item \textbf{Pasillo Central ($A_1$):} Es la ruta directa. OptiBot la utiliza el $40\%$ de las veces ($P(A_1) = 0.40$). Como se analizó en el ejemplo anterior, al ser muy concurrida presenta una probabilidad de demora del $40\%$ ($P(B \mid A_1) = 0.40$).
  \item \textbf{Pasillo Norte ($A_2$):} Es una ruta más amplia y despejada. El robot la utiliza también el $40\%$ de las veces ($P(A_2) = 0.40$). En este pasillo la probabilidad de demora es baja, de solo un $10\%$ ($P(B \mid A_2) = 0.10$).
  \item \textbf{Pasillo Sur ($A_3$):} Es una ruta secundaria que se emplea el $20\%$ restante de las veces ($P(A_3) = 0.20$). Su probabilidad de demora es intermedia, del $20\%$ ($P(B \mid A_3) = 0.20$).
\end{itemize}

Para hallar la probabilidad total de que un viaje sufra demoras, sin importar qué pasillo haya elegido el robot, se aplica la Ley de Probabilidad Total:
\begin{align*}
  P(B) &= P(B \mid A_1) \cdot P(A_1) + P(B \mid A_2) \cdot P(A_2) + P(B \mid A_3) \cdot P(A_3) \\
  &= (0.40 \times 0.40) + (0.10 \times 0.40) + (0.20 \times 0.20) \\
  &= 0.16 + 0.04 + 0.04 \\
  &= \mathbf{0.24 \quad (24\%)}
\end{align*}

\textbf{Conclusión didáctica:} En promedio, el $24\%$ de todos los viajes en el almacén sufrirán demoras. Esta visión global le permite al algoritmo comprender el rendimiento general de su entorno antes de comenzar a optimizar qué porcentaje de viajes le conviene destinar a cada pasillo.

\begin{ideaclave}
Esta suma ponderada no es una curiosidad abstracta. Representa el núcleo matemático sobre el cual se construyen los algoritmos de decisión en Inteligencia Artificial. En el Capítulo~\ref{cap:mdp}, se comprobará que la célebre \textbf{Ecuación de Bellman} aplica exactamente esta misma lógica para calcular el valor de cualquier estado futuro.
\end{ideaclave}

\section{El Concepto Más Importante: El Valor Esperado}
Si una persona compra un boleto de lotería y gana un premio pequeño, cometería una equivocación grave si concluyera que comprar lotería todos los días es una inversión financiera rentable. Juzgar una estrategia basándose únicamente en el resultado de un intento aislado suele conducir a decisiones erróneas.

En el Aprendizaje por Refuerzo, a un robot no le interesa ganar una recompensa por pura casualidad en una sola ocasión; lo que busca es encontrar una estrategia que produzca los mejores resultados posibles \textbf{en promedio a largo plazo}. Para medir este rendimiento sostenido se utiliza el \textbf{valor esperado} (denotado matemáticamente como $\mathbb{E}[X]$):
\begin{equation}
  \mathbb{E}[X] \doteq \sum_{x} x \cdot P(X = x)
\end{equation}
El símbolo $\doteq$ (un signo de igual con un punto encima) es una convención muy común en la literatura de Aprendizaje por Refuerzo y significa simplemente \textbf{«se define como»}. Indica que la fórmula no es el resultado de un despeje matemático previo, sino la propia definición fundamental del concepto.

El valor esperado no pretende adivinar qué valor exacto se observará en el siguiente viaje; lo que indica con precisión es el promedio al que tenderán los resultados si la acción se repite cientos o miles de veces bajo condiciones similares.

\subsection*{Ejemplo: Cálculo del valor esperado en los envíos de OptiBot}
Para comprender cómo se aplica esta fórmula en la toma de decisiones, considérese el balance económico diario de OptiBot al transportar pedidos hacia los camiones de reparto. Se presentan tres escenarios posibles en cada viaje:
\begin{itemize}
  \item \textbf{Entrega rápida (10 minutos):} el cliente abona una tarifa preferencial y deja una ganancia neta de $+\$10$. Esto ocurre el $26\%$ de las veces ($P = 0.26$).
  \item \textbf{Entrega regular (20 minutos):} el viaje se completa dentro del horario estándar y deja una ganancia de $+\$5$. Esto ocurre el $50\%$ de las veces ($P = 0.50$).
  \item \textbf{Entrega con demora (más de 30 minutos):} el almacén debe compensar al cliente con un cupón, resultando en una pérdida de $-\$2$. Esto ocurre el $24\%$ de las veces ($P = 0.24$).
\end{itemize}

Para calcular cuánta ganancia genera OptiBot en promedio por cada paquete transportado, se multiplica cada resultado económico por su respectiva probabilidad y se suman:
\begin{align*}
  \mathbb{E}[\text{Ganancia}] &= (10 \times 0.26) + (5 \times 0.50) + (-2 \times 0.24) \\
  &= 2.60 + 2.50 - 0.48 \\
  &= \mathbf{\$4.62 \text{ por paquete}}
\end{align*}

\textbf{Conclusión didáctica:} Obsérvese un detalle fundamental: \textbf{ningún paquete individual produce exactamente $\$4.62$}. Cada envío aislado deja una ganancia de $\$10$, de $\$5$ o una pérdida de $-\$2$. Sin embargo, si OptiBot transporta 1,000 paquetes a lo largo del mes, la empresa habrá acumulado aproximadamente $\$4,620$. Esta cifra promedio es la que le permite al algoritmo decidir objetivamente si una ruta o una política de trabajo es beneficiosa.

\section{De la Teoría al Código: Ley de los Grandes Números}
Una duda natural es: ¿realmente se cumple este promedio teórico cuando el sistema opera en la práctica?

La \textbf{Ley de los Grandes Números} responde afirmativamente: a medida que un experimento se repite una gran cantidad de veces, la media observada en la realidad se aproxima de manera casi perfecta al valor esperado teórico calculado en el papel.

El siguiente programa en Python simula el trabajo de OptiBot realizando 50,000 despachos en el almacén para verificar experimentalmente si el promedio empírico coincide con los $\$4.62$ calculados a mano:

\begin{lstlisting}[style=rlcode]
import numpy as np

# 1. Definicion de los resultados posibles y sus probabilidades
ganancias = np.array([10.0, 5.0, -2.0])
probabilidades = np.array([0.26, 0.50, 0.24])

# 2. Computo del valor esperado teorico
E_teorico = np.sum(ganancias * probabilidades)
print(f"Valor Esperado Teorico: ${E_teorico:.2f}")

# 3. Simulacion de 50,000 despachos realizados por OptiBot
np.random.seed(42)
despachos = np.random.choice(ganancias, size=50000, p=probabilidades)

# 4. Promedio real obtenido a partir de la simulacion
promedio_observado = np.mean(despachos)
print(f"Promedio Real Observado (50,000 entregas): ${promedio_observado:.2f}")

# 5. Comparacion entre la prediccion teorica y el resultado practico
diferencia = abs(E_teorico - promedio_observado)
print(f"Diferencia absoluta: ${diferencia:.4f}")
\end{lstlisting}

Al ejecutar la simulación, se comprueba que el promedio observado difiere del valor teórico por menos de dos centavos de dólar. Si agrega un cero al número de iteraciones de la simulación, la diferencia será mucho menor y el valor teórico sera prácticamente igual al valor observado.

\begin{notabox}
Con esta base conceptual, el azar deja de ser un obstáculo desconcertante y se transforma en una variable cuantificable y predecible a gran escala. En el Capítulo~\ref{cap:bandidos}, estas herramientas permitirán analizar el primer problema clásico del aprendizaje por refuerzo: cómo debe actuar un agente inteligente ante el dilema entre \textbf{explotar} una opción que ya sabe que funciona o \textbf{explorar} nuevas alternativas para descubrir si ofrecen ganancias superiores.
\end{notabox}

\chapter{El Problema de los Bandidos Multi-Brazo}
\label{cap:bandidos}

\section{La Esencia del Aprendizaje por Refuerzo}
A diferencia del aprendizaje supervisado, donde un maestro externo proporciona la etiqueta correcta para cada entrada, en el Aprendizaje por Refuerzo el agente recibe una señal \textbf{evaluativa} (una recompensa numérica), pero nunca \textbf{instructiva} (nadie le dice cuál habría sido la mejor acción posible).

Para aislar esta dinámica de aprendizaje sin la complicación añadida del movimiento espacial, Sutton y Barto introducen el problema del \textbf{Bandido de $k$-brazos} (\textit{$k$-armed bandit}). En este escenario existe un único estado estático: el agente debe elegir repetidamente una acción entre $k$ opciones disponibles, recibiendo una recompensa numérica extraída de una distribución de probabilidad fija pero desconocida.

\begin{optibot}
OptiBot se encuentra instalado en la estación central de despacho. Por cada paquete que sale de la cinta, debe seleccionar una de \textbf{4 empresas de mensajería}:
\begin{itemize}
  \item \textbf{Proveedor A:} Tiempo rápido, pero cobra recargos ocasionales ($q_*(A) = 3.2$).
  \item \textbf{Proveedor B:} El más eficiente y rentable en promedio ($q_*(B) = 4.8$).
  \item \textbf{Proveedor C:} Económico pero con frecuentes averías y retrasos ($q_*(C) = 2.1$).
  \item \textbf{Proveedor D:} Servicio premium constante ($q_*(D) = 4.1$).
\end{itemize}
OptiBot \textbf{no conoce} estos valores teóricos $q_*$. Cada vez que selecciona un proveedor, observa la ganancia neta obtenida $R_t$. Su meta es maximizar la ganancia total acumulada a lo largo de 1,000 despachos.
\end{optibot}

\section{Valor de Acción y Estimación de Recompensas}
Denotamos la acción elegida en el instante de tiempo $t$ como $A_t$ y la recompensa resultante como $R_t$. El \textbf{valor verdadero} de una acción arbitraria $a$, denotado $q_*(a)$, se define formalmente como la esperanza condicionada:
\begin{equation}
  q_*(a) \doteq \mathbb{E}\left[ R_t \mid A_t = a \right]
\end{equation}
Dado que el agente desconoce $q_*(a)$, mantiene una estimación en el paso $t$, denotada como $Q_t(a)$. Si el agente selecciona una acción $a$ un total de $N_t(a)$ veces hasta el tiempo $t$, la estimación más natural es el promedio de las recompensas recibidas:
\begin{equation}
  Q_t(a) \doteq \frac{\sum_{i=1}^{t-1} R_i \cdot \mathbb{I}(A_i = a)}{N_t(a)}
\end{equation}

\section{La Regla de Actualización Incremental}
Almacenar todas las recompensas históricas en un vector requeriría una memoria que crece indefinidamente con el tiempo. Afortunadamente, podemos actualizar el promedio de forma recursiva:
\begin{align}
  Q_{n+1} &= \frac{1}{n} \sum_{i=1}^n R_i \notag \\
  &= \frac{1}{n} \left( R_n + \sum_{i=1}^{n-1} R_i \right) \notag \\
  &= \frac{1}{n} \left( R_n + (n-1) Q_n \right) \notag \\
  &= Q_n + \frac{1}{n} \left( R_n - Q_n \right)
\end{align}

\begin{ideaclave}
La ecuación anterior representa el patrón estructural más importante de todo el Aprendizaje por Refuerzo:
\begin{equation}
  \text{NuevaEstimación} \leftarrow \text{EstimaciónPrevia} + \text{TamañoDePaso} \times \left[ \text{Objetivo} - \text{EstimaciónPrevia} \right]
\end{equation}
El término entre corchetes, $[\text{Objetivo} - \text{EstimaciónPrevia}]$, se denomina el \textbf{error de estimación}. Este mismo esquema gobierna la Diferencia Temporal, SARSA, Q-Learning y el entrenamiento de redes neuronales profundas.
\end{ideaclave}

Si el entorno no es estacionario (las probabilidades de los proveedores cambian con el tiempo), reemplazamos $\frac{1}{n}$ por una tasa constante $\alpha \in (0, 1]$, lo que confiere mayor peso exponencial a las observaciones más recientes.

\section{El Dilema Exploración vs. Explotación: Métodos $\varepsilon$-Greedy}
\begin{itemize}
  \item \textbf{Explotación ($\text{greedy}$):} Seleccionar la acción con mayor valor estimado actual:
  \begin{equation}
    A_t^* = \arg\max_a Q_t(a)
  \end{equation}
  Aprovecha el conocimiento adquirido para obtener recompensas inmediatas elevadas, pero corre el riesgo de ignorar una acción superior aún no descubierta.
  \item \textbf{Exploración:} Seleccionar deliberadamente una acción subóptima según los datos actuales para recopilar información y mejorar las estimaciones futuras.
\end{itemize}

El algoritmo $\varepsilon$-greedy equilibra ambas necesidades de forma sencilla y elegante:
\begin{equation}
  A_t = \begin{cases}
    \text{acción aleatoria uniforme en } \mathcal{A}, & \text{con probabilidad } \varepsilon, \\
    \arg\max_a Q_t(a), & \text{con probabilidad } 1 - \varepsilon.
  \end{cases}
\end{equation}

\begin{calculomanual}
Supongamos que OptiBot tiene la tabla de valores $Q = [0.0, 0.0, 0.0, 0.0]$ y elige el Proveedor B ($a=2$).
\begin{enumerate}
  \item Observa una recompensa $R_1 = 5.2$. Como $N(B)=1$:
  \[
    Q_2(B) = 0.0 + \frac{1}{1}(5.2 - 0.0) = \mathbf{5.20}
  \]
  \item En el siguiente paso, elige de nuevo el Proveedor B y obtiene $R_2 = 4.4$. Ahora $N(B)=2$:
  \[
    Q_3(B) = 5.20 + \frac{1}{2}(4.4 - 5.20) = 5.20 - 0.40 = \mathbf{4.80}
  \]
  \item En el tercer paso explora el Proveedor A ($a=1$) y obtiene $R_3 = 3.0$. Ahora $N(A)=1$:
  \[
    Q_4(A) = 0.0 + \frac{1}{1}(3.0 - 0.0) = \mathbf{3.00}
  \]
\end{enumerate}
La tabla actualizada es $Q = [3.00, 4.80, 0.00, 0.00]$. El agente explota correctamente a B con probabilidad $1-\varepsilon$, pero mantiene la capacidad de probar C y D.
\end{calculomanual}

\section{Implementación en Python}
A continuación implementamos el agente de bandidos multi-brazo completo en Python:

\begin{lstlisting}[style=rlcode]
import numpy as np
import matplotlib.pyplot as plt

class BanditOptiBot:
    def __init__(self, k=4, epsilon=0.1, alpha=None):
        self.k = k
        self.epsilon = epsilon
        self.alpha = alpha
        self.q_true = np.array([3.2, 4.8, 2.1, 4.1]) # Valores reales ocultos
        self.Q = np.zeros(k)
        self.N = np.zeros(k)

    def select_action(self):
        if np.random.rand() < self.epsilon:
            return np.random.randint(self.k)
        return np.argmax(self.Q)

    def step(self, action):
        # Recompensa estocastica con ruido gaussiano
        reward = np.random.normal(loc=self.q_true[action], scale=1.0)
        self.N[action] += 1
        step_size = self.alpha if self.alpha is not None else (1.0 / self.N[action])
        self.Q[action] += step_size * (reward - self.Q[action])
        return reward

# Simulacion de 2,000 despachos
bot = BanditOptiBot(k=4, epsilon=0.10)
rewards = []
for t in range(2000):
    a = bot.select_action()
    r = bot.step(a)
    rewards.append(r)

print("Estimaciones finales Q:", np.round(bot.Q, 3))
print("Veces seleccionado cada proveedor:", bot.N)
\end{lstlisting}

\begin{notabox}
\textbf{La limitación que motiva el siguiente capítulo:} En el problema de los bandidos, el agente toma decisiones, pero \textbf{el mundo nunca cambia de estado}. Todas las decisiones ocurren en el mismo vacío estático. En la vida real, los sistemas evolucionan: los equipos se desgastan y el entorno transita por diferentes estados físicos. Ese fenómeno dinámico es el objeto de estudio de las \textbf{Cadenas de Markov}.
\end{notabox}

\chapter{Cadenas de Markov y Procesos Estocásticos}
\label{cap:markov}

\section{Introducción a la Dinámica Temporal}
En el capítulo anterior comprobamos que OptiBot podía aprender a seleccionar el mejor proveedor de mensajería mediante la regla de actualización incremental. No obstante, el entorno era completamente estático: cada decisión se tomaba en el mismo instante abstracto y no afectaba a los intentos venideros.

En el mundo físico, los entornos poseen memoria física y estados cambiantes: la batería se descarga, las ruedas pierden adherencia y las cintas transportadoras sufren desgaste mecánico. Para modelar matemáticamente cómo un sistema físico evoluciona en el tiempo bajo la influencia del azar, recurrimos a las \textbf{Cadenas de Markov}.

\section{La Propiedad de Markov}
Consideremos una secuencia temporal de estados discretos $S_0, S_1, S_2, \dots, S_t, S_{t+1}$. Decimos que un proceso estocástico satisface la \textbf{Propiedad de Markov} si la distribución de probabilidad del estado futuro depende exclusivamente del estado presente, siendo condicionalmente independiente de toda la historia pasada:
\begin{equation}
  P\left( S_{t+1} = s_{t+1} \;\middle|\; S_t = s_t, \, S_{t-1} = s_{t-1}, \, \dots, \, S_0 = s_0 \right) = P\left( S_{t+1} = s_{t+1} \;\middle|\; S_t = s_t \right)
\end{equation}

\begin{ideaclave}
La propiedad de Markov no afirma que el pasado carezca de importancia. Lo que establece es que \textbf{toda la información relevante del pasado ya se encuentra encapsulada en el estado presente}. Si conoces con exactitud dónde estás hoy, saber cómo llegaste allí no añade ningún poder predictivo sobre a dónde irás mañana.
\end{ideaclave}

\begin{optibot}
OptiBot supervisa los rodamientos de la cinta transportadora central. En cada turno de trabajo, el estado del rodamiento se clasifica en una de tres categorías discretas $\mathcal{S} = \{ \text{Nuevo } (N), \text{Desgastado } (D), \text{Averiado } (A) \}$.
\begin{itemize}
  \item Si un rodamiento está Nuevo, tiene una probabilidad de $0.70$ de seguir Nuevo en el próximo turno, $0.20$ de pasar a Desgastado y $0.10$ de sufrir una rotura imprevista (Averiado).
  \item Si está Desgastado, tiene probabilidad $0.60$ de seguir Desgastado, $0.30$ de romperse (Averiado) y $0.10$ de que un operario le aplique lubricante y vuelva a considerarse Nuevo.
  \item Si está Averiado, la cinta se detiene hasta que el equipo de mantenimiento lo repara; permanece averiado con probabilidad $0.90$ y es sustituido por uno Nuevo con probabilidad $0.10$.
\end{itemize}
\end{optibot}

\section{La Matriz de Transición de Estados}
Para un conjunto finito de estados $\mathcal{S} = \{s_1, s_2, \dots, s_n\}$, organizamos todas las probabilidades de transición condicionales en una matriz cuadrada $P \in \mathbb{R}^{n \times n}$, donde cada elemento representa:
\begin{equation}
  P_{ij} \doteq P(S_{t+1} = s_j \mid S_t = s_i)
\end{equation}
Dado que partiendo de cualquier estado $s_i$ el sistema debe transitar obligatoriamente a algún estado del espacio muestral, cada fila de $P$ satisface la condición estocástica:
\begin{equation}
  \sum_{j=1}^n P_{ij} = 1 \quad \forall i \in \{1, \dots, n\}
\end{equation}

Para el caso de la cinta transportadora de OptiBot, la matriz de transición es:
\begin{equation}
  P = \begin{pmatrix}
    0.70 & 0.20 & 0.10 \\
    0.10 & 0.60 & 0.30 \\
    0.10 & 0.00 & 0.90
  \end{pmatrix}
\end{equation}

\section{Transiciones a Múltiples Pasos y Distribución Estacionaria}
Si denotamos la distribución de probabilidad inicial del sistema como un vector fila $\mu_0 = [P(S_0=s_1), \dots, P(S_0=s_n)]$, la distribución en el paso 1 viene dada por el producto vectorial-matricial:
\begin{equation}
  \mu_1 = \mu_0 P
\end{equation}
Por inducción y por las ecuaciones de Chapman-Kolmogorov, la distribución en el paso $k$ es:
\begin{equation}
  \mu_k = \mu_0 P^k
\end{equation}

Cuando la cadena de Markov es irreducible y aperiódica, existe una única \textbf{distribución estacionaria} $\pi^* = [\pi_1^*, \dots, \pi_n^*]$ que es invariante bajo la aplicación de la matriz de transición:
\begin{equation}
  \pi^* P = \pi^* \quad \text{sujeto a} \quad \sum_{i=1}^n \pi_i^* = 1
\end{equation}

\begin{calculomanual}
Supongamos que hoy instalamos un rodamiento totalmente nuevo, por lo que el vector de estado inicial es $\mu_0 = [1.0, 0.0, 0.0]$. ¿Cuál es la probabilidad exacta de que el rodamiento esté averiado dentro de 2 turnos?
\begin{align*}
  \mu_1 &= \mu_0 P = [0.70, 0.20, 0.10] \\
  \mu_2 &= \mu_1 P = [0.70, 0.20, 0.10] \begin{pmatrix}
    0.70 & 0.20 & 0.10 \\
    0.10 & 0.60 & 0.30 \\
    0.10 & 0.00 & 0.90
  \end{pmatrix}
\end{align*}
Calculamos la tercera componente de $\mu_2$, correspondiente al estado Averiado ($A$):
\begin{align*}
  P(S_2 = A) &= (0.70)(0.10) + (0.20)(0.30) + (0.10)(0.90) \\
  &= 0.07 + 0.06 + 0.09 = \mathbf{0.22}
\end{align*}
Existe un $22\%$ de probabilidad de que la cinta esté averiada en el segundo turno.
\end{calculomanual}

\section{Implementación en Python}
A continuación calculamos la evolución a largo plazo y la distribución estacionaria mediante autovalores y autovectores en NumPy:

\begin{lstlisting}[style=rlcode]
import numpy as np

P = np.array([
    [0.70, 0.20, 0.10],
    [0.10, 0.60, 0.30],
    [0.10, 0.00, 0.90]
])

# 1. Proyeccion a 2 pasos
mu_0 = np.array([1.0, 0.0, 0.0])
P_2 = np.linalg.matrix_power(P, 2)
mu_2 = mu_0 @ P_2
print("Distribucion en t=2:", np.round(mu_2, 3))

# 2. Calculo de la distribucion estacionaria (pi * P = pi)
# Equivale a (P.T - I) * pi.T = 0
evals, evecs = np.linalg.eig(P.T)
# Buscamos el autovector correspondiente al autovalor 1
idx = np.argmin(np.abs(evals - 1.0))
pi_star = np.real(evecs[:, idx])
pi_star = pi_star / np.sum(pi_star) # Normalizar para que sume 1

print("Distribucion Estacionaria pi*:", np.round(pi_star, 4))
print(f"Porcentaje de tiempo averiada a largo plazo: {pi_star[2]*100:.1f}%")
\end{lstlisting}

\begin{notabox}
\textbf{La limitación que motiva el siguiente paso:} En este capítulo hemos dotado al mundo de dinamismo temporal. Sabemos cómo evoluciona el estado del entorno probabilísticamente. Sin embargo, \textbf{somos meros espectadores pasivos}: OptiBot no puede intervenir para reparar la cinta ni para cambiar su trayectoria. ¿Qué ocurre si unimos las \textbf{acciones con recompensas} del Capítulo~\ref{cap:bandidos} con los \textbf{estados dinámicos} del Capítulo~\ref{cap:markov}? La respuesta es la estructura nuclear del aprendizaje por refuerzo: el \textbf{Proceso de Decisión de Markov}.
\end{notabox}

\chapter{Procesos de Decisión de Markov y Ecuaciones de Bellman}
\label{cap:mdp}

\section{La Unificación de Estados, Acciones y Recompensas}
Los capítulos precedentes han desarrollado de forma independiente dos dimensiones complementarias de la toma de decisiones:
\begin{itemize}
  \item En el Capítulo~\ref{cap:bandidos}, se abordó cómo elegir entre diferentes \textbf{acciones} evaluando beneficios numéricos inmediatos, pero en un entorno sin movimiento ni \textbf{estados}.
  \item En el Capítulo~\ref{cap:markov}, se analizó la evolución temporal de \textbf{estados} físicos en el almacén, pero en un régimen puramente pasivo donde el sistema cambiaba sin que un agente pudiera intervenir con \textbf{acciones}.
\end{itemize}

El \textbf{Proceso de Decisión de Markov} (\textit{Markov Decision Process}, o MDP) representa la síntesis definitiva de todo el aprendizaje: en cada instante de tiempo, el agente observa en qué \textbf{estado} se encuentra, evalúa las alternativas disponibles, elige una \textbf{acción} deliberada, el entorno cambia hacia un nuevo estado y devuelve una señal evaluativa formal llamada \textbf{recompensa}.

\subsection{La Hipótesis de la Recompensa}
¿Cómo se le comunica a un robot o algoritmo qué es lo que se espera que logre? En el Aprendizaje por Refuerzo, todo objetivo se formaliza a través de una señal escalar denominada \textbf{recompensa} ($R_t$).

Este principio se conoce como la \textbf{Hipótesis de la Recompensa}:
\begin{quote}
  \textit{«Todo lo que entendemos por metas y propósitos puede formalizarse adecuadamente como la maximización del valor esperado de la suma acumulada de una señal numérica escalar (la recompensa).»}
\end{quote}
La recompensa no le dice al agente qué pasos específicos debe ejecutar (no es una instrucción); únicamente le comunica qué tan deseable es el resultado alcanzado. Con esta señal, el agente debe descubrir la secuencia de decisiones óptima.

\begin{optibot}
En esta fase del desarrollo, OptiBot opera como un robot móvil autónomo encargado de completar el \textbf{ciclo de despacho de pedidos} a lo largo de cuatro zonas operativas del almacén:
\begin{enumerate}
  \item \textbf{Muelle de Carga ($S_1$):} Zona inicial de recepción. OptiBot debe decidir cómo cargar y despachar el pedido.
  \item \textbf{Pasillo Central ($S_2$):} Zona de navegación y transporte. OptiBot debe seleccionar el modo de conducción para trasladar el paquete.
  \item \textbf{Bahía de Entrega ($S_3$):} Área final de descarga del producto antes de salir hacia el cliente.
  \item \textbf{Fin de Turno / Misión Cumplida ($S_4$):} Estado terminal absorbente donde la entrega se ha completado exitosamente y concluye el episodio ($V(S_4) = 0$).
\end{enumerate}
En cada uno de los estados activos ($S_1, S_2, S_3$), OptiBot dispone de \textbf{dos acciones alternativas} que alteran sus probabilidades de avance y su consumo energético.
\end{optibot}

\section{Definición Formal de un MDP: La Tupla de Siete Elementos}
A diferencia de las Cadenas de Markov, que modelan sistemas pasivos únicamente mediante estados y transiciones, un Proceso de Decisión de Markov permite formalizar la toma activa de decisiones mediante una tupla de siete elementos:
\begin{equation}
  \mathcal{M} = \langle \mathcal{S}, \, \mathcal{A}, \, \mathcal{P}, \, \mathcal{R}, \, \gamma, \, \mu, \, \pi \rangle
\end{equation}
cuyos componentes se definen como:
\begin{itemize}
  \item $\mathcal{S}$ (\textbf{Espacio de Estados}): Conjunto de situaciones o posiciones en las cuales puede encontrarse el sistema ($\mathcal{S} = \{S_1, S_2, S_3, S_4\}$).
  \item $\mathcal{A}$ (\textbf{Espacio de Acciones}): Conjunto de decisiones de control ejecutables por el agente en cada estado.
  \item $\mathcal{P}$: Función de dinámica de transición que describe la probabilidad condicional de que el sistema pase al estado sucesor $s'$ habiendo tomado la acción $a$ en el estado actual $s$:
  \begin{equation}
    P(s' \mid s, a) \doteq P(S_{t+1} = s' \mid S_t = s, A_t = a)
  \end{equation}
  \item $\mathcal{R}$: Función de recompensa que asigna un valor escalar esperado a cada transición ejecutada bajo el par estado-acción:
  \begin{equation}
    R(s, a, s') \doteq \mathbb{E}[R_{t+1} \mid S_t = s, A_t = a, S_{t+1} = s']
  \end{equation}
  \item $\gamma \in [0, 1)$ (\textbf{Factor de Descuento}): Ponderador temporal que regula la importancia relativa de las recompensas futuras frente a las inmediatas.
  \item $\mu$: Distribución inicial de estados, que asigna la probabilidad de comenzar en cada estado del sistema (en OptiBot: $\mu = [1, 0, 0, 0]$, pues siempre parte del Muelle).
  \item $\pi$: Política del agente, que define su estrategia asignando a cada estado una distribución de probabilidad sobre las acciones elegibles: $\pi(a \mid s) \doteq P(A_t = a \mid S_t = s)$.
\end{itemize}

\subsection{El Grafo del MDP: Estados, Acciones y Recompensas}
La incorporación de acciones transforma la simple cadena de Markov en un grafo de decisión. La Figura~\ref{fig:optibot_mdp_graph} ilustra el modelo de OptiBot:
\begin{itemize}
  \item Los \textbf{círculos continuos} representan los estados activos ($S_1, S_2, S_3$), y el círculo doble representa el estado terminal ($S_4$).
  \item Las \textbf{flechas continuas} indican las acciones disponibles en cada estado.
  \item Las \textbf{flechas discontinuas} representan las transiciones estocásticas del entorno hacia el estado sucesor, rotuladas con el par $(p, \, r)$: la probabilidad de transición $p(s' \mid s, a)$ y la recompensa escalar obtenida $r$.
\end{itemize}

\begin{figure}[!htb]
\centering
\begin{tikzpicture}[
  scale=0.72,
  transform shape,
  >=Stealth,
  every node/.style={font=\sffamily},
  stateNode/.style={circle, draw=TealDark, fill=Teal!12, very thick, minimum size=1.8cm, align=center, font=\small\bfseries\sffamily, text=Ink},
  termNode/.style={circle, draw=Coral, fill=Coral!15, double, very thick, minimum size=1.8cm, align=center, font=\small\bfseries\sffamily, text=Ink},
  actNode/.style={rectangle, draw=Coral, fill=Coral!12, rounded corners=3pt, font=\scriptsize\bfseries\sffamily, text=Coral!90!black, inner sep=5pt},
  solidEdge/.style={->, thick, color=Coral},
  dashEdge/.style={->, thick, dashed, color=TealDark},
  lbl/.style={fill=white, draw=TealDark!40, thin, inner sep=3pt, rounded corners=2pt, font=\scriptsize\sffamily\bfseries, text=Ink}
]

  % --- ESTADOS ---
  \node[stateNode] (S1) at (0, 0)       {Muelle\\($S_1$)};
  \node[stateNode] (S2) at (9.0, 5.5)   {Pasillo\\($S_2$)};
  \node[stateNode] (S3) at (9.0, -5.5)  {Entrega\\($S_3$)};
  \node[termNode]  (S4) at (17.5, -5.5) {Fin\\($S_4$)};

  % --- SALIDAS DE S1 ---
  \node[actNode] (a1_std) at (4.0, 3.6)  {Carga Estándar};
  \node[actNode] (a1_rap) at (4.0, -3.6) {Línea Rápida};
  \draw[solidEdge] (S1) -- (a1_std);
  \draw[solidEdge] (S1) -- (a1_rap);
  \draw[dashEdge]  (a1_std) -- node[above=4pt, lbl] {$1.00,\; +1.0$} (S2.195);
  \draw[dashEdge]  (a1_rap) -- node[below=4pt, lbl] {$1.00,\; +0.1$} (S3.165);

  % --- ACCIONES DE S2 (Pasillo) ---
  % Carril Oeste: Modo Eco
  \node[actNode] (a2_eco) at (5.2, 1.0) {Modo Eco};
  \draw[solidEdge] (S2.215) to[bend right=15] (a2_eco.north);
  \draw[dashEdge]  (a2_eco.east) to[bend right=15] node[right=4pt, lbl] {$0.20,\; +0.3$} (S2.245);
  \draw[dashEdge]  (a2_eco.south) to[bend right=15] node[left=4pt, lbl] {$0.80,\; +1.0$} (S3.140);

  % Columna Central: Modo Normal
  \node[actNode] (a2_nor) at (9.0, 0.0) {Modo Normal};
  \draw[solidEdge] (S2.270) -- (a2_nor.north);
  \draw[dashEdge]  (a2_nor.east) to[out=50, in=-60] node[right=4pt, lbl] {$0.70,\; +0.5$} (S2.300);
  \draw[dashEdge]  (a2_nor.south) -- node[fill=white, lbl] {$0.30,\; +0.3$} (S3.90);

  % --- ACCIONES DE S3 (Entrega) ---
  % Carril Este Interior: Descarga Manual
  \node[actNode] (a3_man) at (13.2, -3.2) {Descarga Manual};
  \draw[solidEdge] (S3.40) to[bend left=15] (a3_man.west);
  \draw[dashEdge]  (a3_man.east) to[bend left=12] node[above=3pt, lbl] {$0.85,\; +1.0$} (S4.120);
  \draw[dashEdge]  (a3_man.north) to[out=95, in=-25] node[pos=0.35, right=4pt, lbl] {$0.15,\; -1.0$} (S2.330);

  % Carril Este Exterior: Descarga Auto (rodea S4 por fuera)
  \node[actNode] (a3_aut) at (13.2, -7.8) {Descarga Auto};
  \draw[solidEdge] (S3.320) to[bend right=15] (a3_aut.west);
  \draw[dashEdge]  (a3_aut.east) to[bend right=12] node[below=3pt, lbl] {$0.90,\; +1.0$} (S4.240);
  \draw[dashEdge]  (a3_aut.south east) to[out=-10, in=-90] (20.2, -5.5) to[out=90, in=5] node[pos=0.45, right=4pt, lbl] {$0.10,\; -1.0$} (S2.10);

\end{tikzpicture}
\caption{Grafo de transición del Proceso de Decisión de Markov para el ciclo logístico de OptiBot. Las flechas continuas conectan cada estado con sus acciones disponibles. Las flechas discontinuas representan las probabilidades de transición y recompensas $(p, \, r)$ resultantes hacia los estados sucesores.}
\label{fig:optibot_mdp_graph}
\end{figure}

\subsection{Representación Matricial del MDP}
El MDP de la Figura~\ref{fig:optibot_mdp_graph} se puede representar de forma matricial definiendo una \textbf{matriz de transición} ($MT_a$) y una \textbf{matriz de recompensa} ($R_a$) para cada una de las acciones disponibles:

\subsubsection*{1. Matrices de Transición por Acción ($MT_a$)}
Cada celda $(i, j)$ indica la probabilidad de saltar del estado $S_i$ al estado $S_j$ si el agente ejecuta dicha acción:
\begin{align*}
  MT_{\text{Estándar}} &= \begin{pmatrix}
    0 & 1.00 & 0 & 0 \\
    0 & 0 & 0 & 0 \\
    0 & 0 & 0 & 0 \\
    0 & 0 & 0 & 0
  \end{pmatrix}, &
  MT_{\text{Rápida}} &= \begin{pmatrix}
    0 & 0 & 1.00 & 0 \\
    0 & 0 & 0 & 0 \\
    0 & 0 & 0 & 0 \\
    0 & 0 & 0 & 0
  \end{pmatrix} \\[6pt]
  MT_{\text{Eco}} &= \begin{pmatrix}
    0 & 0 & 0 & 0 \\
    0 & 0.20 & 0.80 & 0 \\
    0 & 0 & 0 & 0 \\
    0 & 0 & 0 & 0
  \end{pmatrix}, &
  MT_{\text{Normal}} &= \begin{pmatrix}
    0 & 0 & 0 & 0 \\
    0 & 0.70 & 0.30 & 0 \\
    0 & 0 & 0 & 0 \\
    0 & 0 & 0 & 0
  \end{pmatrix} \\[6pt]
  MT_{\text{Manual}} &= \begin{pmatrix}
    0 & 0 & 0 & 0 \\
    0 & 0 & 0 & 0 \\
    0 & 0.15 & 0 & 0.85 \\
    0 & 0 & 0 & 0
  \end{pmatrix}, &
  MT_{\text{Auto}} &= \begin{pmatrix}
    0 & 0 & 0 & 0 \\
    0 & 0 & 0 & 0 \\
    0 & 0.10 & 0 & 0.90 \\
    0 & 0 & 0 & 0
  \end{pmatrix}
\end{align*}

\subsubsection*{2. Matrices de Recompensa por Acción ($R_a$)}
De forma análoga, cada celda $(i, j)$ define la recompensa obtenida al transitar de $S_i$ a $S_j$ bajo esa acción:
\begin{align*}
  R_{\text{Estándar}} &= \begin{pmatrix}
    0 & 1.00 & 0 & 0 \\
    0 & 0 & 0 & 0 \\
    0 & 0 & 0 & 0 \\
    0 & 0 & 0 & 0
  \end{pmatrix}, &
  R_{\text{Rápida}} &= \begin{pmatrix}
    0 & 0 & 0.10 & 0 \\
    0 & 0 & 0 & 0 \\
    0 & 0 & 0 & 0 \\
    0 & 0 & 0 & 0
  \end{pmatrix} \\[6pt]
  R_{\text{Eco}} &= \begin{pmatrix}
    0 & 0 & 0 & 0 \\
    0 & 0.30 & 1.00 & 0 \\
    0 & 0 & 0 & 0 \\
    0 & 0 & 0 & 0
  \end{pmatrix}, &
  R_{\text{Normal}} &= \begin{pmatrix}
    0 & 0 & 0 & 0 \\
    0 & 0.50 & 0.30 & 0 \\
    0 & 0 & 0 & 0 \\
    0 & 0 & 0 & 0
  \end{pmatrix} \\[6pt]
  R_{\text{Manual}} &= \begin{pmatrix}
    0 & 0 & 0 & 0 \\
    0 & 0 & 0 & 0 \\
    0 & -1.00 & 0 & 1.00 \\
    0 & 0 & 0 & 0
  \end{pmatrix}, &
  R_{\text{Auto}} &= \begin{pmatrix}
    0 & 0 & 0 & 0 \\
    0 & 0 & 0 & 0 \\
    0 & -1.00 & 0 & 1.00 \\
    0 & 0 & 0 & 0
  \end{pmatrix}
\end{align*}

\subsection{Tabla de Parámetros del MDP}
El Cuadro~\ref{tab:mdp_parameters} resume todas las dinámicas y recompensas del entorno en una estructura tabular compacta:


\begin{table}[!htb]
\centering
\small
\begin{tabular}{llccl}
\toprule
\textbf{Estado Inicial ($s$)} & \textbf{Acción ($a$)} & \textbf{Probabilidad ($p$)} & \textbf{Recompensa ($r$)} & \textbf{Nuevo Estado ($s'$)} \\
\midrule
Muelle ($S_1$) & Carga Estándar & $1.00$ & $+1.00$ & Pasillo ($S_2$) \\
Muelle ($S_1$) & Línea Rápida   & $1.00$ & $+0.10$ & Entrega ($S_3$) \\
\midrule
Pasillo ($S_2$) & Modo Eco    & $0.20$ & $+0.30$ & Pasillo ($S_2$) \\
Pasillo ($S_2$) & Modo Eco    & $0.80$ & $+1.00$ & Entrega ($S_3$) \\
Pasillo ($S_2$) & Modo Normal & $0.70$ & $+0.50$ & Pasillo ($S_2$) \\
Pasillo ($S_2$) & Modo Normal & $0.30$ & $+0.30$ & Entrega ($S_3$) \\
\midrule
Entrega ($S_3$) & Descarga Manual & $0.15$ & $-1.00$ & Pasillo ($S_2$) \\
Entrega ($S_3$) & Descarga Manual & $0.85$ & $+1.00$ & Fin ($S_4$) \\
Entrega ($S_3$) & Descarga Auto   & $0.10$ & $-1.00$ & Pasillo ($S_2$) \\
Entrega ($S_3$) & Descarga Auto   & $0.90$ & $+1.00$ & Fin ($S_4$) \\
\bottomrule
\end{tabular}
\caption{Parámetros de transición y recompensas de cada acción del MDP de OptiBot.}
\label{tab:mdp_parameters}
\end{table}

\subsection{El Retorno Descontado y la Preferencia Temporal}
El objetivo de OptiBot no es obtener un premio aislado, sino maximizar el \textbf{retorno acumulado descontado} $G_t$:
\begin{equation}
  G_t \doteq R_{t+1} + \gamma R_{t+2} + \gamma^2 R_{t+3} + \dots = \sum_{k=0}^{\infty} \gamma^k R_{t+k+1}
\end{equation}
El parámetro $\gamma \in [0, 1)$ pondera la preferencia temporal: valores cercanos a $0$ generan agentes miopes centrados en la gratificación inmediata, mientras que valores próximos a $1$ promueven agentes previsores y pacientes. Para este problema adoptamos $\gamma = 0.90$.

\section{Políticas y Funciones de Valor}
Una \textbf{política} $\pi$ formaliza el comportamiento del agente. Para evaluar y comparar políticas, se introducen dos funciones de valor fundamentales:
\begin{enumerate}
  \item \textbf{Función de Valor de Estado $V^\pi(s)$:} Valor esperado del retorno al comenzar en el estado $s$ y seguir la política $\pi$:
  \begin{equation}
    V^\pi(s) \doteq \mathbb{E}_\pi [G_t \mid S_t = s]
  \end{equation}
  \item \textbf{Función de Valor de Acción $Q^\pi(s, a)$:} Valor esperado del retorno al ejecutar la acción $a$ en el estado $s$, siguiendo $\pi$ en adelante:
  \begin{equation}
    Q^\pi(s, a) \doteq \mathbb{E}_\pi [G_t \mid S_t = s, A_t = a]
  \end{equation}
\end{enumerate}

\section{La Ecuación de Expectativa de Bellman}
La genialidad de Richard Bellman consistió en demostrar que el valor de un estado no requiere simular infinitos pasos futuros: se expresa como una relación recursiva de dos términos:
\begin{equation}
  V^\pi(s) = \sum_{a \in \mathcal{A}} \pi(a \mid s) \sum_{s', r} p(s', r \mid s, a) \Big[ r + \gamma V^\pi(s') \Big]
\end{equation}
Es decir: el valor de estar hoy en $s$ es igual a la recompensa inmediata promedio que se recibe al dar el paso ($r$), más el valor futuro descontado de dónde aterrice el agente ($\gamma V(s')$).

\section{Evaluación Iterativa de Políticas y Determinación de la Política Óptima}
Para resolver el MDP y descubrir la política óptima $\pi^*$ que indique qué acción debe ejecutarse en cada estado, se emplea el algoritmo clásico de \textbf{Evaluación Iterativa de Políticas / Iteración de Valores} formulado por Sutton y Barto:

\begin{tcolorbox}[
  colback=white,
  colframe=Ink,
  arc=3pt,
  boxrule=1pt,
  left=12pt, right=12pt, top=8pt, bottom=8pt,
  title={\textbf{\textcolor{white}{Evaluación Iterativa de Políticas, para estimar $V \approx v_\pi$}}},
  fonttitle=\sffamily\bfseries,
  coltitle=white
]
\textbf{Entrada:} $\pi$, la política a ser evaluada.\\
\textbf{Parámetro del algoritmo:} un umbral pequeño $\theta > 0$ que determina la precisión de la estimación.\\
\textbf{Inicializar:} $V(s)$, para todo $s \in \mathcal{S}^+$, arbitrariamente excepto que $V(\text{terminal}) = 0$.\\[6pt]
\textbf{Bucle:}\\[2pt]
\hspace*{1.5em}$\Delta \leftarrow 0$\\[2pt]
\hspace*{1.5em}\textbf{Bucle para cada} $s \in \mathcal{S}$:\\[2pt]
\hspace*{3.0em}$v \leftarrow V(s)$\\[2pt]
\hspace*{3.0em}$V(s) \leftarrow \max_{a} \sum_{s', r} p(s', r \mid s, a) \Big[ r + \gamma V(s') \Big]$\\[2pt]
\hspace*{3.0em}$\Delta \leftarrow \max\big(\Delta, \, |v - V(s)|\big)$\\[4pt]
\textbf{hasta que} $\Delta < \theta$
\end{tcolorbox}

\subsection{Resolución Manual Paso a Paso: Aplicación en OptiBot}
Consideramos un umbral de convergencia $\theta = 0.01$ y un factor de descuento $\gamma = 0.90$.

\paragraph{Paso 1: Inicialización.}
Como primer paso, se establece la función valor $V(s)$ de cada estado en algún valor arbitrario. Por simplicidad, inicializamos todos los valores en cero con la condición estricta de que el estado terminal siempre vale cero:
\begin{table}[!htb]
\centering
\begin{tabular}{lc}
\toprule
\textbf{Estado ($s$)} & \textbf{$V_0(s)$} \\
\midrule
Muelle ($S_1$) & $0.00$ \\
Pasillo ($S_2$) & $0.00$ \\
Entrega ($S_3$) & $0.00$ \\
Fin ($S_4$) & $0.00$ (Terminal) \\
\bottomrule
\end{tabular}
\caption{Función de valor inicial de cada estado del MDP de OptiBot.}
\label{tab:v0_init}
\end{table}

\paragraph{Paso 2: Primera Iteración ($k = 1$).}
En cada iteración, el parámetro de variación $\Delta$ se inicializa en cero. Para cada estado, se calcula el valor de cada acción posible evaluando el término:
\begin{equation}
  p(s', r \mid s, a) \Big[ r + \gamma V_0(s') \Big]
\end{equation}
y se suman los resultados para obtener el rendimiento acumulado $\Sigma$ de dicha acción. Luego, se selecciona el valor más alto ($\max$) entre las acciones disponibles y se asigna como el nuevo valor $V(s)$ de ese estado.

El Cuadro~\ref{tab:iteracion_1} muestra el desglose aritmético completo de la primera iteración:

\begin{table}[!htb]
\centering
\small
\setlength{\tabcolsep}{3.5pt}
\begin{tabular}{lllccccc}
\toprule
$s$ & $a$ & $s'$ & $r$ & $p(s' \mid s, a)$ & $V_0(s')$ & $p(s' \mid s, a)[r + \gamma V_0(s')]$ & $\Sigma$ \\
\midrule
Muelle & Carga Estándar & Pasillo & $1.00$ & $1.00$ & $0.0000$ & $1.000000$ & $\mathbf{1.000000}$ \\
Muelle & Línea Rápida   & Entrega & $0.10$ & $1.00$ & $0.0000$ & $0.100000$ & $0.100000$ \\
\midrule
Pasillo & Modo Eco    & Pasillo & $0.30$ & $0.20$ & $0.0000$ & $0.060000$ & $0.060000$ \\
Pasillo & Modo Eco    & Entrega & $1.00$ & $0.80$ & $0.0000$ & $0.800000$ & $\mathbf{0.860000}$ \\
Pasillo & Modo Normal & Pasillo & $0.50$ & $0.70$ & $0.0000$ & $0.350000$ & $0.350000$ \\
Pasillo & Modo Normal & Entrega & $0.30$ & $0.30$ & $0.0000$ & $0.090000$ & $0.440000$ \\
\midrule
Entrega & Descarga Manual & Pasillo & $-1.00$ & $0.15$ & $0.0000$ & $-0.150000$ & $-0.150000$ \\
Entrega & Descarga Manual & Fin     & $1.00$  & $0.85$ & $0.0000$ & $0.850000$  & $0.700000$ \\
Entrega & Descarga Auto   & Pasillo & $-1.00$ & $0.10$ & $0.0000$ & $-0.100000$ & $-0.100000$ \\
Entrega & Descarga Auto   & Fin     & $1.00$  & $0.90$ & $0.0000$ & $0.900000$  & $\mathbf{0.800000}$ \\
\bottomrule
\end{tabular}
\vspace{6pt}

\textbf{Actualización de valores de $V(s)$:}

\begin{tabular}{lcccc}
\toprule
$s$ & $V(s)_{\text{anterior}}$ & $V(s)_{\text{nuevo}}$ & $\text{Diferencia}$ & $\Delta = \max(\Delta, \text{diferencia})$ \\
\midrule
Muelle ($S_1$)  & $0.000000$ & $1.000000$ & $1.000000$ & $1.000000$ \\
Pasillo ($S_2$) & $0.000000$ & $0.860000$ & $0.860000$ & $1.000000$ \\
Entrega ($S_3$) & $0.000000$ & $0.800000$ & $0.800000$ & $1.000000$ \\
\bottomrule
\end{tabular}
\caption{Primera iteración del algoritmo de evaluación iterativa para OptiBot.}
\label{tab:iteracion_1}
\end{table}

Como $\Delta = 1.0000 \ge 0.01$, el algoritmo no se detiene y avanza a la siguiente iteración utilizando el nuevo vector de valores $[1.0000, \, 0.8600, \, 0.8000]$.

\paragraph{Paso 3: Iteración Final y Convergencia.}
El ciclo de actualización se repite de forma progresiva. En cada pasada, los valores futuros descubiertos se propagan hacia atrás a través del factor de descuento $\gamma = 0.90$. Tras 7 iteraciones, la variación máxima se reduce a $\Delta = 0.009071$, cumpliendo con éxito la condición de parada $\Delta < 0.01$.

El Cuadro~\ref{tab:iteracion_final} documenta la iteración final del algoritmo:

\begin{table}[!htb]
\centering
\small
\setlength{\tabcolsep}{3.5pt}
\begin{tabular}{lllccccc}
\toprule
$s$ & $a$ & $s'$ & $r$ & $p(s' \mid s, a)$ & $V_6(s')$ & $p(s' \mid s, a)[r + \gamma V_6(s')]$ & $\Sigma$ \\
\midrule
Muelle & Carga Estándar & Pasillo & $1.00$ & $1.00$ & $1.895925$ & $2.706333$ & $\mathbf{2.706333}$ \\
Muelle & Línea Rápida   & Entrega & $0.10$ & $1.00$ & $0.969695$ & $0.972726$ & $0.972726$ \\
\midrule
Pasillo & Modo Eco    & Pasillo & $0.30$ & $0.20$ & $1.895925$ & $0.401267$ & $0.401267$ \\
Pasillo & Modo Eco    & Entrega & $1.00$ & $0.80$ & $0.969695$ & $1.498180$ & $\mathbf{1.899447}$ \\
Pasillo & Modo Normal & Pasillo & $0.50$ & $0.70$ & $1.895925$ & $1.544433$ & $1.544433$ \\
Pasillo & Modo Normal & Entrega & $0.30$ & $0.30$ & $0.969695$ & $0.351818$ & $1.896251$ \\
\midrule
Entrega & Descarga Manual & Pasillo & $-1.00$ & $0.15$ & $1.895925$ & $0.105950$ & $0.105950$ \\
Entrega & Descarga Manual & Fin     & $1.00$  & $0.85$ & $0.000000$ & $0.850000$ & $0.955950$ \\
Entrega & Descarga Auto   & Pasillo & $-1.00$ & $0.10$ & $1.895925$ & $0.070633$ & $0.070633$ \\
Entrega & Descarga Auto   & Fin     & $1.00$  & $0.90$ & $0.000000$ & $0.900000$ & $\mathbf{0.970633}$ \\
\bottomrule
\end{tabular}
\vspace{6pt}

\textbf{Actualización de valores de $V(s)$:}

\begin{tabular}{lcccc}
\toprule
$s$ & $V(s)_{\text{anterior}}$ & $V(s)_{\text{nuevo}}$ & $\text{Diferencia}$ & $\Delta = \max(\Delta, \text{diferencia})$ \\
\midrule
Muelle ($S_1$)  & $2.697262$ & $2.706333$ & $0.009071$ & $0.009071$ \\
Pasillo ($S_2$) & $1.895925$ & $1.899447$ & $0.003522$ & $0.009071$ \\
Entrega ($S_3$) & $0.969695$ & $0.970633$ & $0.000938$ & $0.009071$ \\
\bottomrule
\end{tabular}
\caption{Iteración final (Iteración 7) del algoritmo con convergencia garantizada ($\Delta = 0.009071 < 0.01$).}
\label{tab:iteracion_final}
\end{table}

\subsection{Deducción de la Política Óptima}
Una vez que el algoritmo converge, la política óptima $\pi^*$ se obtiene identificando la acción que obtuvo el valor más alto en la columna $\Sigma$ para cada estado del sistema:
\begin{equation}
  \pi^*(s) = \begin{cases}
    (Muelle, \, \text{Carga Estándar}) & \text{con } V^*(S_1) \approx 2.71 \\
    (Pasillo, \, \text{Modo Eco})      & \text{con } V^*(S_2) \approx 1.90 \\
    (Entrega, \, \text{Descarga Automática}) & \text{con } V^*(S_3) \approx 0.97
  \end{cases}
\end{equation}

\subsection*{Interpretación Operativa de la Política Óptima}
A diferencia de un modelo con incentivos mal calibrados, este resultado tiene una coherencia lógica absoluta en el contexto operativo del almacén:
\begin{enumerate}
  \item \textbf{En el Muelle ($S_1$):} OptiBot prefiere la \textbf{Carga Estándar} sobre la Línea Rápida ($2.71$ vs $0.97$). Aunque la línea rápida lo teletransporta directamente a la entrega, el costo de activación penaliza fuertemente la ganancia inicial ($r = 0.10$ vs $1.00$), haciendo que el despacho normal sea económicamente superior a largo plazo.
  \item \textbf{En el Pasillo ($S_2$):} OptiBot selecciona el \textbf{Modo Eco} frente al Modo Normal ($1.90$ vs $1.89$). A pesar de que el modo normal ofrece una recompensa inmediata mayor por demorarse ($0.50$ vs $0.30$), el Modo Eco ofrece una probabilidad de avance del $80\%$ hacia la bahía de entrega frente a solo el $30\%$ del modo normal. El agente prefiere llegar antes a la meta que demorarse cobrando pequeñas recompensas intermedias.
  \item \textbf{En la Bahía de Entrega ($S_3$):} OptiBot elige la \textbf{Descarga Automática} ($0.97$ vs $0.95$). Al reducir el riesgo de atascos al $10\%$ (frente al $15\%$ del método manual), asegura con un $90\%$ de probabilidad el éxito final del pedido.
  \item \textbf{Evolución decreciente hacia la meta ($2.71 \to 1.90 \to 0.97 \to 0$):} El valor de cada estado representa el \textbf{retorno restante esperado} hasta culminar el turno. Cuanto más cerca está el robot de completar la entrega, menos etapas le quedan por recorrer y menor es el retorno futuro acumulable, alcanzando exactamente cero en el estado terminal Fin.
\end{enumerate}

\section{Implementación en Python}
El siguiente script en Python implementa el algoritmo completo de evaluación iterativa para OptiBot, calcula la convergencia y extrae de forma automática la política óptima $\pi^*$:

\begin{lstlisting}[style=rlcode]
import numpy as np

# 1. Definicion del MDP de OptiBot
estados = ["Muelle", "Pasillo", "Entrega", "Fin"]
num_estados = 3  # Estados activos (Fin es terminal con V=0)
gamma = 0.90
theta = 0.01

# Tabla de transiciones: [origen, accion, destino, recompensa, probabilidad]
transiciones = [
    # Muelle (S1)
    (0, "Carga Estandar", 1,  1.00, 1.00),
    (0, "Linea Rapida",   2,  0.10, 1.00),
    # Pasillo (S2)
    (1, "Modo Eco",       1,  0.30, 0.20),
    (1, "Modo Eco",       2,  1.00, 0.80),
    (1, "Modo Normal",    1,  0.50, 0.70),
    (1, "Modo Normal",    2,  0.30, 0.30),
    # Entrega (S3)
    (2, "Descarga Manual",3,  1.00, 0.85),
    (2, "Descarga Manual",1, -1.00, 0.15),
    (2, "Descarga Auto",  3,  1.00, 0.90),
    (2, "Descarga Auto",  1, -1.00, 0.10),
]

# 2. Inicializacion de funciones de valor: V(s) = 0
V = np.zeros(4) # V[3] corresponde al estado terminal 'Fin' = 0

# 3. Bucle de Evaluacion Iterativa / Value Iteration
iteracion = 0
while True:
    delta = 0.0
    iteracion += 1
    V_nuevo = V.copy()
    
    # Evaluar cada estado activo
    for s in range(num_estados):
        v_anterior = V[s]
        # Identificar acciones disponibles en el estado s
        acciones_s = sorted(list(set(a for s_orig, a, _, _, _ in transiciones if s_orig == s)))
        q_valores = []
        
        for accion in acciones_s:
            # Calcular suma de p * [r + gamma * V(s')]
            suma_q = sum(p * (r + gamma * V[dest]) 
                         for s_orig, a, dest, r, p in transiciones 
                         if s_orig == s and a == accion)
            q_valores.append(suma_q)
            
        V_nuevo[s] = max(q_valores)
        delta = max(delta, abs(v_anterior - V_nuevo[s]))
        
    V = V_nuevo
    if delta < theta:
        break

# 4. Extraccion de la Politica Optima pi*(s)
politica_optima = {}
for s in range(num_estados):
    acciones_s = sorted(list(set(a for s_orig, a, _, _, _ in transiciones if s_orig == s)))
    mejor_q = -np.inf
    mejor_accion = None
    for accion in acciones_s:
        suma_q = sum(p * (r + gamma * V[dest]) 
                     for s_orig, a, dest, r, p in transiciones 
                     if s_orig == s and a == accion)
        if suma_q > mejor_q:
            mejor_q = suma_q
            mejor_accion = accion
    politica_optima[estados[s]] = (mejor_accion, mejor_q)

# 5. Visualizacion de resultados
print(f"=== CONVERGENCIA ALCANZADA EN {iteracion} ITERACIONES (theta = {theta}) ===")
for estado, (accion, valor) in politica_optima.items():
    print(f"  Estado {estado:8s} -> Accion Optima: {accion:18s} | V*(s) = {valor:.4f}")
\end{lstlisting}

\subsection*{Salida del Programa e Interpretación}
Al ejecutar el código, la terminal entrega el siguiente reporte verificado:
\begin{lstlisting}[style=rlcode]
=== CONVERGENCIA ALCANZADA EN 7 ITERACIONES (theta = 0.01) ===
  Estado Muelle   -> Accion Optima: Carga Estandar   | V*(s) = 2.7095
  Estado Pasillo  -> Accion Optima: Modo Eco         | V*(s) = 1.9008
  Estado Entrega  -> Accion Optima: Descarga Auto    | V*(s) = 0.9710
\end{lstlisting}

Los resultados del algoritmo computacional reproducen con precisión exacta los cálculos manuales derivados en las tablas precedentes, confirmando que la política óptima para OptiBot consiste en preparar los pedidos con Carga Estándar, navegar el almacén en Modo Eco y completar la descarga mediante el sistema Automático.

\begin{notabox}
\textbf{El puente hacia el siguiente capítulo:} En este capítulo hemos definido rigurosamente qué es un Proceso de Decisión de Markov y cómo evaluar sus funciones de valor mediante las Ecuaciones de Bellman bajo el supuesto de conocer el modelo analítico del entorno. En el Capítulo~\ref{cap:rl_tradicional}, abordaremos el desafío fundamental del \textbf{Aprendizaje por Refuerzo Tradicional}: ¿cómo puede OptiBot aprender la política óptima sin disponer de planos ni modelos analíticos de probabilidad? Introduciremos el aprendizaje libre de modelo mediante Diferencia Temporal y deduciremos en detalle el algoritmo fundamental \textbf{Q-Learning}.
\end{notabox}

\chapter{Programación Dinámica e Interacción Política Generalizada}
\label{cap:dp}

\section{El Problema del Control Óptimo}
En el capítulo anterior aprendimos a resolver el problema de la \textbf{evaluación de predicción}: dada una política fija $\pi$, calcular su función de valor $V^\pi(s)$. Sin embargo, un agente inteligente autónomo no debe limitarse a resignarse ante una política arbitraria; debe descubrir la mejor política posible.

Decimos que una política $\pi^*$ es \textbf{óptima} si su valor esperado acumulado es mayor o igual que el de cualquier otra política para todos y cada uno de los estados del entorno:
\begin{equation}
  V^{\pi^*}(s) \ge V^\pi(s) \quad \forall s \in \mathcal{S}, \quad \forall \pi
\end{equation}
Todas las políticas óptimas comparten la misma función de valor óptima de estado, denotada $V^*(s) \doteq \max_\pi V^\pi(s)$, y la misma función de valor de acción óptima $Q^*(s,a) \doteq \max_\pi Q^\pi(s,a)$.

\section{Ecuaciones de Bellman de Optimalidad}
Bajo una política óptima, la decisión tomada en cada estado debe coincidir con la acción que maximiza la suma de la recompensa esperada inmediata más el valor óptimo del estado sucesor descontado:
\begin{align}
  V^*(s) &= \max_{a \in \mathcal{A}} Q^*(s, a) \notag \\
  &= \max_{a \in \mathcal{A}} \left[ \mathcal{R}(s, a) + \gamma \sum_{s' \in \mathcal{S}} \mathcal{P}_{ss'}^a V^*(s') \right]
\end{align}
De forma complementaria, para el valor de acción óptimo:
\begin{equation}
  Q^*(s, a) = \mathcal{R}(s, a) + \gamma \sum_{s' \in \mathcal{S}} \mathcal{P}_{ss'}^a \max_{a' \in \mathcal{A}} Q^*(s', a')
\end{equation}

\begin{ideaclave}
Nótese una diferencia matemática crítica: la Ecuación de Expectativa de Bellman era lineal, lo que permitía resolverla mediante inversión matricial. Por el contrario, la Ecuación de Optimalidad contiene el operador no lineal $\max_{a}$, lo que invalida el álgebra lineal directa y nos obliga a recurrir a algoritmos iterativos de punto fijo.
\end{ideaclave}

\section{Policy Iteration y Value Iteration}
La Programación Dinámica (DP) ofrece dos algoritmos fundamentales para hallar $V^*$ y $\pi^*$, asumiendo que conocemos perfectamente el modelo del entorno ($\mathcal{P}$ y $\mathcal{R}$):

\subsection{1. Policy Iteration (Iteración de Políticas)}
Comienza con una política arbitraria $\pi_0$ y repite dos etapas consecutivas:
\begin{enumerate}
  \item \textbf{Evaluación de Política:} Resuelve iterativamente $V^{\pi_k}$ hasta convergencia:
  \[
    V_{m+1}(s) \leftarrow \sum_{a} \pi_k(a \mid s) \left[ \mathcal{R}(s,a) + \gamma \sum_{s'} \mathcal{P}_{ss'}^a V_m(s') \right]
  \]
  \item \textbf{Mejora de Política:} Construye una política estrictamente mejor eligiendo de forma codiciosa (\textit{greedy}) con respecto al valor calculado:
  \[
    \pi_{k+1}(s) \leftarrow \arg\max_{a \in \mathcal{A}} \left[ \mathcal{R}(s,a) + \gamma \sum_{s'} \mathcal{P}_{ss'}^a V^{\pi_k}(s') \right]
  \]
\end{enumerate}
El \textbf{Teorema de Mejora de Política} garantiza que $V^{\pi_{k+1}}(s) \ge V^{\pi_k}(s)$ para todo estado $s$. Al existir un número finito de políticas deterministas en un MDP finito, el algoritmo converge exactamente a $\pi^*$ en un número finito de iteraciones.

\subsection{2. Value Iteration (Iteración de Valores)}
La evaluación completa de la política hasta su estricta convergencia en cada ciclo es computacionalmente costosa. Richard Bellman propuso truncar la evaluación a una sola pasada, incorporando la actualización de optimalidad directamente:
\begin{equation}
  V_{k+1}(s) \leftarrow \max_{a \in \mathcal{A}} \left[ \mathcal{R}(s, a) + \gamma \sum_{s' \in \mathcal{S}} \mathcal{P}_{ss'}^a V_k(s') \right]
\end{equation}
Por el \textbf{Teorema del Punto Fijo de Banach}, el operador de optimalidad de Bellman es una $\gamma$-contracción estricta en la norma del supremo ($\|T(u) - T(v)\|_\infty \le \gamma \|u - v\|_\infty$). Por tanto, partiendo de cualquier vector inicial $V_0$, la sucesión $\{V_k\}$ converge invariablemente a la solución única $V^*$.

\section{Interacción Política Generalizada (GPI)}
Casi todos los métodos de Aprendizaje por Refuerzo pueden conceptualizarse como instancias de \textbf{Interacción Política Generalizada} (\textit{Generalized Policy Iteration}, GPI). En GPI interactúan dos procesos simultáneos:
\begin{itemize}
  \item El proceso de evaluación busca que la función de valor sea consistente con la política actual.
  \item El proceso de mejora busca que la política sea codiciosa respecto a la función de valor actual.
\end{itemize}
Ambos procesos cooperan y compiten dinámicamente hasta que alcanzan un estado de equilibrio en el que la política es codiciosa con respecto a su propia función de valor, lo que sólo ocurre en $(\pi^*, V^*)$.

\begin{optibot}
El fabricante entrega a los ingenieros los planos completos del almacén: una cuadrícula de $3 \times 3$ celdas donde se conocen con absoluta exactitud las probabilidades de movimiento de OptiBot y las penalizaciones por paso ($-1$ en cada movimiento, $+10$ al llegar al punto de entrega en la celda $(2,2)$). Con esta información perfecta, OptiBot no necesita experimentar a ciegas: puede calcular la política de navegación óptima directamente en su procesador central mediante Value Iteration.
\end{optibot}

\section{Implementación en Python: Value Iteration para Cuadrícula}
\begin{lstlisting}[style=rlcode]
import numpy as np

def value_iteration(grid_size=3, gamma=0.9, theta=1e-4):
    V = np.zeros((grid_size, grid_size))
    target = (grid_size - 1, grid_size - 1)
    actions = [(-1, 0), (1, 0), (0, -1), (0, 1)] # Arriba, Abajo, Izq, Der

    while True:
        delta = 0
        for r in range(grid_size):
            for c in range(grid_size):
                if (r, c) == target:
                    continue
                v_old = V[r, c]
                q_candidates = []
                for dr, dc in actions:
                    nr, nc = max(0, min(grid_size-1, r+dr)), max(0, min(grid_size-1, c+dc))
                    reward = 10.0 if (nr, nc) == target else -1.0
                    q_candidates.append(reward + gamma * V[nr, nc])
                V[r, c] = max(q_candidates)
                delta = max(delta, abs(v_old - V[r, c]))
        if delta < theta:
            break
    return V

V_optimo = value_iteration()
print("Mapa de Valores Optimos V*(s):")
print(np.round(V_optimo, 2))
\end{lstlisting}

\begin{notabox}
\textbf{La limitación que motiva el siguiente capítulo:} La Programación Dinámica es conceptualmente perfecta pero padece una grave limitación práctica: \textbf{requiere un modelo completo y perfecto de la física del entorno} ($\mathcal{P}$ y $\mathcal{R}$). En el mundo real, los almacenes cambian, los suelos se desgastan y nadie conoce las probabilidades exactas de transición. Necesitamos algoritmos capaces de aprender \textbf{sin modelo}, basándose exclusivamente en la experiencia empírica. Ese es el reino de los \textbf{Métodos Monte Carlo}.
\end{notabox}

\chapter{Métodos Monte Carlo}
\label{cap:monte_carlo}

\section{El Aprendizaje Libre de Modelo (Model-Free RL)}
Hasta el Capítulo~\ref{cap:dp}, asumimos una condición privilegiada: conocíamos con precisión absoluta las ecuaciones del entorno ($\mathcal{P}$ y $\mathcal{R}$). En casi todas las aplicaciones prácticas de la robótica, la logística o los videojuegos, esta suposición resulta inviable.

Los \textbf{Métodos Monte Carlo} rompen de raíz esta dependencia: permiten al agente estimar las funciones de valor y optimizar políticas basándose \textbf{únicamente en la experiencia empírica} obtenida al interactuar con el entorno real o con un simulador.

\begin{optibot}
OptiBot es enviado a un nuevo centro de distribución cuyo plano y dinámicas de tráfico son completamente desconocidos. Ya no disponemos de matrices de transición precalculadas. En su lugar, dejamos que OptiBot ejecute turnos completos de trabajo (\textbf{episodios}). Durante el turno, el robot registra minuciosamente la secuencia de estados transitados, acciones tomadas y recompensas recibidas. Al final de la jornada laboral, analiza la experiencia acumulada para deducir qué caminos y decisiones resultaron más rentables.
\end{optibot}

\section{Estimación de Valores por Muestreo de Retornos}
Un episodio completo genera una trayectoria finita de interacción:
\begin{equation}
  S_0, A_0, R_1, S_1, A_1, R_2, S_2, \dots, S_{T-1}, A_{T-1}, R_T, S_T
\end{equation}
donde $S_T$ representa el estado terminal (por ejemplo, el paquete ha sido entregado exitosamente o la batería se ha agotado por completo).

A partir de cualquier instante $t$, el retorno descontado observado en esa trayectoria concreta se calcula de manera exacta:
\begin{equation}
  G_t = R_{t+1} + \gamma R_{t+2} + \dots + \gamma^{T-t-1} R_T
\end{equation}

Recordando el principio visto en el Capítulo~\ref{cap:probabilidad}, el valor de un estado $V^\pi(s)$ es la esperanza matemática $\mathbb{E}_\pi[G_t \mid S_t = s]$. Por la Ley de los Grandes Números, el promedio aritmético de los retornos empíricos observados tras visitar el estado $s$ converge a su valor verdadero a medida que el número de episodios tiende a infinito:
\begin{equation}
  V(s) \approx \frac{1}{N(s)} \sum_{i=1}^{N(s)} G_i(s)
\end{equation}

\subsection{First-Visit vs. Every-Visit Monte Carlo}
\begin{itemize}
  \item \textbf{First-Visit MC:} Para cada episodio, sólo se promedia el retorno $G_t$ correspondiente a la \textbf{primera vez} que el estado $s$ es visitado dentro de dicha trayectoria. Garantiza estimaciones insesgadas e independientes de $V^\pi(s)$.
  \item \textbf{Every-Visit MC:} Se promedian todos los retornos observados en cada visita a $s$, incluso si el agente visita el mismo estado repetidamente en un bucle durante el mismo episodio.
\end{itemize}

\section{Control Monte Carlo sin Modelo: Estimar $Q(s,a)$}
En Programación Dinámica podíamos limitarnos a calcular $V(s)$ porque, para elegir la mejor acción, disponíamos de la matriz $\mathcal{P}_{ss'}^a$ y podíamos calcular:
\[
  a^* = \arg\max_a \sum_{s'} \mathcal{P}_{ss'}^a V(s')
\]
\begin{ideaclave}
En el aprendizaje libre de modelo, no conocemos $\mathcal{P}_{ss'}^a$. Conocer únicamente $V(s)$ es inútil para la toma de decisiones, porque no sabemos a qué estado sucesor nos conducirá cada acción. Por ello, los métodos libres de modelo deben estimar obligatoriamente la \textbf{función de valor de acción $Q(s,a)$}.
\end{ideaclave}

Teniendo $Q(s,a)$, el agente puede actuar de forma codiciosa directamente sin necesidad de consultar ningún modelo físico:
\begin{equation}
  \pi(s) = \arg\max_{a \in \mathcal{A}} Q(s,a)
\end{equation}

Para garantizar que todas las combinaciones de estado y acción sigan siendo visitadas y sus valores continúen refinándose, el agente emplea una política $\varepsilon$-greedy con respecto a la tabla $Q(s,a)$ actual.

\begin{calculomanual}
Supongamos que OptiBot parte del Muelle ($S_0$), toma la acción Avanzar ($A_0$) y genera el siguiente episodio con factor de descuento $\gamma = 0.9$:
\[
  S_0 \xrightarrow{R_1=-1} S_1 \xrightarrow{R_2=-1} S_2 \xrightarrow{R_3=+10} \text{Meta}
\]
Calculamos los retornos acumulados recorriendo la trayectoria \textbf{de atrás hacia adelante}:
\begin{align*}
  G_2 &= R_3 = \mathbf{+10.0} \\
  G_1 &= R_2 + \gamma G_2 = -1.0 + (0.9)(10.0) = \mathbf{+8.0} \\
  G_0 &= R_1 + \gamma G_1 = -1.0 + (0.9)(8.0) = \mathbf{+6.2}
\end{align*}
El par $(S_0, A_0)$ añade la muestra $G_0 = 6.2$ a su registro de retornos. Si antes tenía un promedio $Q(S_0, A_0) = 4.0$ con 2 visitas previas ($N=2$), su nuevo valor estimado aplicando la regla incremental es:
\[
  Q_{\text{nuevo}}(S_0, A_0) = 4.0 + \frac{1}{3}(6.2 - 4.0) = 4.0 + 0.733 = \mathbf{4.733}
\]
\end{calculomanual}

\section{Implementación en Python: Control Monte Carlo First-Visit}
\begin{lstlisting}[style=rlcode]
import numpy as np
from collections import defaultdict

class MonteCarloControl:
    def __init__(self, n_actions=4, gamma=0.9, epsilon=0.15):
        self.n_actions = n_actions
        self.gamma = gamma
        self.epsilon = epsilon
        self.Q = defaultdict(lambda: np.zeros(n_actions))
        self.returns_sum = defaultdict(float)
        self.returns_count = defaultdict(int)

    def select_action(self, state):
        if np.random.rand() < self.epsilon:
            return np.random.randint(self.n_actions)
        return np.argmax(self.Q[state])

    def update_episode(self, episode):
        # episode es una lista de tuplas (estado, accion, recompensa)
        G = 0.0
        visited_sa = set()
        # Recorrido inverso para computar G_t en O(T)
        for state, action, reward in reversed(episode):
            G = self.gamma * G + reward
            sa_pair = (state, action)
            if sa_pair not in visited_sa:
                visited_sa.add(sa_pair)
                self.returns_sum[sa_pair] += G
                self.returns_count[sa_pair] += 1
                # Actualizacion del promedio
                self.Q[state][action] = (
                    self.returns_sum[sa_pair] / self.returns_count[sa_pair]
                )
\end{lstlisting}

\begin{notabox}
\textbf{La limitación que motiva el siguiente capítulo:} Los métodos Monte Carlo son insesgados y fáciles de entender, pero sufren dos debilidades severas:
\begin{enumerate}
  \item \textbf{Alta varianza:} Al sumar muchas recompensas estocásticas a lo largo de todo un episodio, el retorno $G_t$ presenta una varianza elevada, lo que ralentiza la convergencia.
  \item \textbf{Exclusivamente episódicos:} El robot debe esperar obligatoriamente a que concluya el turno completo para actualizar el primer paso. Si un turno dura horas o la tarea es continua y nunca termina, Monte Carlo no puede aprender en tiempo real.
\end{enumerate}
¿Es posible aprender paso a paso, tras cada movimiento individual, sin esperar al final del episodio? La respuesta es el mayor hito conceptual del RL: el \textbf{Aprendizaje por Diferencia Temporal (TD)}.
\end{notabox}

\chapter{Diferencia Temporal: SARSA y Q-Learning}
\label{cap:td}

\section{La Gran Idea: El Aprendizaje por Bootstrapping}
El Aprendizaje por Diferencia Temporal (\textit{Temporal-Difference Learning}, TD) representa la idea más distintiva y profunda del Aprendizaje por Refuerzo. Descrita por Richard Sutton en 1988, la diferencia temporal fusiona magistralmente las mejores virtudes de los dos mundos estudiados en los capítulos anteriores:
\begin{itemize}
  \item Al igual que los \textbf{Métodos Monte Carlo}, TD aprende directamente a partir de la experiencia bruta en línea, sin necesidad de conocer el modelo de transición del entorno.
  \item Al igual que la \textbf{Programación Dinámica}, TD actualiza sus estimaciones basándose en parte en otras estimaciones previas, sin esperar al desenlace final del episodio. Esta propiedad se conoce formalmente como \textbf{bootstrapping}.
\end{itemize}

\begin{optibot}
OptiBot ya no espera a que termine su jornada laboral de 8 horas para calcular los retornos y corregir sus decisiones. Ahora, tras cada metro recorrido y cada segundo transcurrido, observa inmediatamente la recompensa obtenida $R_{t+1}$ y el nuevo estado $S_{t+1}$. Si el nuevo estado parece más ventajoso de lo previsto, incrementa al instante el valor del estado anterior; si detecta peligro o demora, reduce su estimación sin dilación.
\end{optibot}

\section{Predicción con TD(0)}
Mientras que Monte Carlo espera al retorno real completo $G_t$ para realizar la actualización:
\[
  V(S_t) \leftarrow V(S_t) + \alpha \left[ G_t - V(S_t) \right]
\]
el método más simple de Diferencia Temporal, denominado \textbf{TD(0)}, sustituye el retorno futuro desconocido por una aproximación en un solo paso:
\begin{equation}
  V(S_t) \leftarrow V(S_t) + \alpha \left[ \underbrace{R_{t+1} + \gamma V(S_{t+1})}_{\text{Objetivo TD}} - V(S_t) \right]
\end{equation}

La cantidad entre corchetes se denomina el \textbf{Error de Diferencia Temporal}, comúnmente abreviado como \textbf{Error TD} ($\delta_t$):
\begin{equation}
  \delta_t \doteq R_{t+1} + \gamma V(S_{t+1}) - V(S_t)
\end{equation}
El error TD cuantifica la discrepancia entre la estimación que el agente tenía en el paso $t$ y la evidencia mejorada obtenida tras avanzar un paso hacia $t+1$.

\section{Algoritmos de Control: SARSA vs. Q-Learning}
Para controlar el movimiento de OptiBot sin un modelo físico del almacén, debemos estimar las funciones de valor de acción $Q(s, a)$. Existen dos enfoques fundamentales según la relación entre la política que genera los datos y la política que se optimiza:

\subsection{1. SARSA: Control On-Policy}
El nombre SARSA proviene de la tupla de 5 elementos que intervienen en cada actualización: $(S_t, A_t, R_{t+1}, S_{t+1}, A_{t+1})$. La regla de actualización es:
\begin{equation}
  Q(S_t, A_t) \leftarrow Q(S_t, A_t) + \alpha \left[ R_{t+1} + \gamma Q(S_{t+1}, A_{t+1}) - Q(S_t, A_t) \right]
\end{equation}
En SARSA, la siguiente acción $A_{t+1}$ se selecciona utilizando la \textbf{misma política de comportamiento} $\varepsilon$-greedy con la que el agente está actuando. Por ello, SARSA aprende el valor de la política que realmente ejecuta, internalizando el riesgo de que el robot pueda cometer un error exploratorio aleatorio. En entornos con precipicios o riesgos graves, SARSA aprende rutas más \textbf{prudentes y conservadoras}.

\subsection{2. Q-Learning: Control Off-Policy}
Propuesto por Christopher Watkins en 1989, \textbf{Q-Learning} es uno de los algoritmos más influyentes de la inteligencia artificial. Su regla de actualización es:
\begin{equation}
  Q(S_t, A_t) \leftarrow Q(S_t, A_t) + \alpha \left[ R_{t+1} + \gamma \max_{a' \in \mathcal{A}} Q(S_{t+1}, a') - Q(S_t, A_t) \right]
\end{equation}

\begin{ideaclave}
Nótese la diferencia fundamental entre ambos algoritmos:
\begin{itemize}
  \item En SARSA, el objetivo utiliza $Q(S_{t+1}, A_{t+1})$, donde $A_{t+1}$ es la acción realmente seleccionada (que puede ser un movimiento exploratorio aleatorio y subóptimo).
  \item En Q-Learning, el objetivo utiliza $\max_{a'} Q(S_{t+1}, a')$, asumiendo que en el siguiente estado el agente actuará de forma estrictamente codiciosa (\textit{greedy}).
\end{itemize}
Esto permite a Q-Learning aprender la \textbf{política óptima directa $Q^*$} de forma \textit{off-policy}, con independencia de los desvíos exploratorios que el robot cometa durante el entrenamiento.
\end{ideaclave}

\begin{calculomanual}
Supongamos que OptiBot se encuentra en el Pasillo A ($S=0$) y ejecuta la acción Avanzar ($A=1$). Recibe una recompensa inmediata $R = -1.0$ y llega al Pasillo B ($S'=1$).
\begin{itemize}
  \item Valor actual en la tabla: $Q(0, 1) = 3.50$.
  \item En el Pasillo B, las estimaciones para sus 3 acciones son:
  \[
    Q(1, \text{GirarIzquierda}) = 6.00, \quad Q(1, \text{Avanzar}) = 8.00, \quad Q(1, \text{GirarDerecha}) = 4.00
  \]
  \item Parámetros de aprendizaje: tasa de paso $\alpha = 0.20$, factor de descuento $\gamma = 0.90$.
\end{itemize}
Desarrollamos el paso de actualización de Q-Learning con lápiz y papel:
\begin{enumerate}
  \item Calculamos el mejor valor de acción en el estado siguiente:
  \[
    \max_{a'} Q(S', a') = \max(6.00, \, 8.00, \, 4.00) = \mathbf{8.00}
  \]
  \item Construimos el objetivo de Bellman TD:
  \[
    \text{Objetivo} = R + \gamma \max_{a'} Q(S', a') = -1.0 + (0.90)(8.00) = -1.0 + 7.20 = \mathbf{6.20}
  \]
  \item Calculamos el error TD de sorpresa:
  \[
    \delta = \text{Objetivo} - Q(S, A) = 6.20 - 3.50 = \mathbf{+2.70}
  \]
  \item Actualizamos la celda correspondiente en la tabla Q:
  \[
    Q_{\text{nuevo}}(0, 1) = Q(0, 1) + \alpha \cdot \delta = 3.50 + (0.20)(+2.70) = 3.50 + 0.54 = \mathbf{4.04}
  \]
\end{enumerate}
La estimación del robot sube de $3.50$ a $4.04$, reflejando que el Pasillo B ofrece opciones de navegación altamente ventajosas.
\end{calculomanual}

\section{Implementación en Python: Q-Learning Tabular}
\begin{lstlisting}[style=rlcode]
import numpy as np

class QLearningAgent:
    def __init__(self, n_states, n_actions, alpha=0.1, gamma=0.95, epsilon=0.1):
        self.n_states = n_states
        self.n_actions = n_actions
        self.alpha = alpha
        self.gamma = gamma
        self.epsilon = epsilon
        self.Q = np.zeros((n_states, n_actions))

    def select_action(self, state):
        if np.random.rand() < self.epsilon:
            return np.random.randint(self.n_actions)
        return np.argmax(self.Q[state])

    def update(self, s, a, r, s_next, done):
        best_next_q = 0.0 if done else np.max(self.Q[s_next])
        td_target = r + self.gamma * best_next_q
        td_error = td_target - self.Q[s, a]
        self.Q[s, a] += self.alpha * td_error

# Bucle principal de interaccion con el entorno de OptiBot
# env = WarehouseEnv()
# for ep in range(1000):
#     s, _ = env.reset()
#     done = False
#     while not done:
#         a = agent.select_action(s)
#         s_next, r, done, _, _ = env.step(a)
#         agent.update(s, a, r, s_next, done)
#         s = s_next
\end{lstlisting}

\begin{notabox}
\textbf{La limitación que motiva el siguiente capítulo:} Q-Learning tabular es extraordinariamente robusto para problemas discretos pequeños. Sin embargo, su talón de Aquiles es la \textbf{maldición de la dimensionalidad}. Si el almacén de OptiBot mide $100 \times 100$ metros y el robot añade sensores continuos de distancia LiDAR, orientación continua en radianes $\theta \in [0, 2\pi)$ y lecturas de nivel de batería en tanto por ciento ($0.0\%$ a $100.0\%$), el número de estados posibles es literalmente \textbf{infinito}. Una tabla matricial no cabe en ninguna memoria RAM y jamás visitaría cada celda. Necesitamos pasar del almacenamiento tabular a la \textbf{Aproximación de Funciones con Redes Neuronales Profundas (DQN)}.
\end{notabox}

\chapter{Deep Q-Networks: De Tablas a Redes Neuronales}
\label{cap:dqn}

\section{La Maldición de la Dimensionalidad y la Necesidad de Generalizar}
En el Capítulo~\ref{cap:td} consolidamos el algoritmo Q-Learning tabular. Asignar un número real a cada par estado-acción en una matriz bidimensional es intuitivo y convergente, pero resulta inviable en la robótica moderna:
\begin{itemize}
  \item Si un robot percibe una imagen de cámara de $84 \times 84$ píxeles en escala de grises con 256 niveles de brillo, el número de estados posibles es $256^{84 \times 84} \approx 10^{16973}$, una cifra enormemente superior al número de átomos en el universo observable.
  \item Si el robot mide magnitudes físicas continuas (como la velocidad en m/s o el ángulo de giro en grados), el espacio de estados es incontable: la probabilidad de visitar exactamente el mismo estado real dos veces es prácticamente cero.
\end{itemize}

La solución consiste en sustituir la tabla rígida por una \textbf{función paramétrica diferenciable}:
\begin{equation}
  Q(s, a; \boldsymbol{\theta}) \approx Q^*(s, a)
\end{equation}
donde $\boldsymbol{\theta}$ representa el vector de pesos sinápticos de una red neuronal artificial profunda. La ventaja colosal de las redes neuronales es su capacidad de \textbf{generalización inductiva}: la red aprende patrones geométricos y conceptuales, infiriendo valores $Q$ precisos para estados que jamás presenció durante el entrenamiento basándose en su similitud con experiencias anteriores.

\begin{optibot}
OptiBot recibe una actualización completa de hardware. Se eliminan las celdas discretas de la cuadrícula y se equipa con:
\begin{enumerate}
  \item Un sensor LiDAR frontal continuo que mide la distancia al obstáculo más cercano en metros ($d \in [0.05, 20.0]$ m).
  \item Un sensor de batería que reporta el porcentaje continuo exacto de carga ($b \in [0.0\%, 100.0\%]$).
  \item Un encoder de odometría que registra la velocidad lineal actual ($v \in [0.0, 3.0]$ m/s).
  \item Una brújula digital que calcula el ángulo relativo hacia el muelle de descarga ($\theta \in [-\pi, \pi]$ rad).
\end{enumerate}
El estado es un vector continuo cuatridimensional $\mathbf{s} = [d, b, v, \theta]^T \in \mathbb{R}^4$. Las acciones siguen siendo discretas: $\{ \text{Avanzar}, \text{Girar Izquierda}, \text{Girar Derecha}, \text{Recargar} \}$.
\end{optibot}

\section{La Tríada Mortal de Sutton \& Barto}
Durante más de dos décadas, entrenar redes neuronales profundas con algoritmos de aprendizaje por refuerzo fue considerado impracticable debido a divergencias numéricas e inestabilidad extrema. Richard Sutton y Andrew Barto identificaron la causa teórica subyacente, conocida como la \textbf{Tríada Mortal}. La divergencia matemática acecha siempre que coinciden simultáneamente estos tres factores:
\begin{enumerate}
  \item \textbf{Aproximación de funciones no lineal:} El uso de redes neuronales en lugar de tablas discretas exactas.
  \item \textbf{Bootstrapping:} Utilizar la predicción del propio modelo en $t+1$ como objetivo para corregir la predicción en $t$ (la esencia de Bellman y TD).
  \item \textbf{Aprendizaje fuera de política (\textit{Off-Policy}):} Aprender sobre la política óptima codiciosa a partir de datos generados por una política exploratoria diferente.
\end{enumerate}

\section{Los Dos Pilares Estabilizadores de DQN (DeepMind 2015)}
En 2015, Volodymyr Mnih y el equipo de Google DeepMind publicaron en la revista \textit{Nature} el algoritmo \textbf{Deep Q-Network (DQN)}, demostrando por primera vez que una red neuronal profunda podía aprender a jugar decenas de videojuegos de Atari al nivel de un experto humano directamente desde los píxeles crudos. El éxito de DQN radica en dos mecanismos esenciales que desactivan la Tríada Mortal:

\subsection{1. Buffer de Repetición de Experiencias (\textit{Experience Replay})}
En el aprendizaje supervisado estándar, las muestras del conjunto de entrenamiento son independientes e idénticamente distribuidas (i.i.d.). En RL, por el contrario, los datos llegan en secuencias temporales fuertemente correlacionadas: el estado en $t+1$ depende críticamente de lo que el robot hizo en $t$. Entrenar una red neuronal con datos secuenciales provoca \textbf{olvido catastrófico} y oscilaciones destructivas.

El \textbf{Replay Buffer} $\mathcal{D}$ almacena las últimas $N$ transiciones observadas por el robot:
\[
  e_t = (s_t, a_t, r_{t+1}, s_{t+1}, \text{done}_{t+1})
\]
Durante el entrenamiento, en lugar de actualizar los pesos con la transición inmediata, el algoritmo extrae un \textbf{mini-batch aleatorio uniforme} de $\mathcal{D}$. Esto rompe por completo las correlaciones temporales y permite reutilizar cada experiencia valiosa en múltiples actualizaciones de gradiente.

\subsection{2. Red Objetivo Congelada (\textit{Target Network})}
En el Q-Learning tabular, modificar la celda $Q(s,a)$ no alteraba el valor numérico de las demás celdas. Sin embargo, en una red neuronal, un paso de retropropagación que actualiza los pesos $\boldsymbol{\theta}$ para corregir $Q(s, a)$ altera sutilmente las salidas de la red para \textbf{todos los demás estados y acciones}. Dado que el objetivo TD depende de la propia red ($\max_{a'} Q(s', a'; \boldsymbol{\theta})$), actualizar $\boldsymbol{\theta}$ equivale a ``perseguir un blanco móvil que se aleja con cada disparo''.

DQN desacopla el objetivo manteniendo dos redes idénticas:
\begin{itemize}
  \item La \textbf{Red Online} ($\boldsymbol{\theta}$): Se actualiza en cada paso mediante descenso de gradiente.
  \item La \textbf{Red Objetivo o Target Network} ($\boldsymbol{\theta}^-$): Sus pesos permanecen fijos y congelados durante $C$ pasos. Se utiliza exclusivamente para calcular el valor del objetivo TD:
  \begin{equation}
    y_i = r + \gamma (1 - \text{done}) \max_{a'} Q(s', a'; \boldsymbol{\theta}^-)
  \end{equation}
\end{itemize}
Cada $C$ pasos temporales (por ejemplo, cada 1,000 pasos), los pesos de la red online se copian a la red objetivo: $\boldsymbol{\theta}^- \leftarrow \boldsymbol{\theta}$.

\section{Función de Pérdida de Bellman y Double DQN}
La función de pérdida que se minimiza mediante descenso de gradiente es el \textbf{Error Cuadrático Medio de Bellman}:
\begin{equation}
  \mathcal{L}(\boldsymbol{\theta}) = \mathbb{E}_{(s, a, r, s', d) \sim \mathcal{D}} \left[ \Big( r + \gamma (1 - d) \max_{a'} Q(s', a'; \boldsymbol{\theta}^-) - Q(s, a; \boldsymbol{\theta}) \Big)^2 \right]
\end{equation}

\subsection{Double DQN (Hado van Hasselt, 2016)}
El operador $\max_{a'}$ en el objetivo de DQN induce un \textbf{sesgo de sobreestimación} sistemático: si el valor Q de una acción contiene ruido estocástico positivo, el operador de máximo lo seleccionará preferentemente, inflando artificialmente el valor estimado del futuro.

\textbf{Double DQN} elimina este sesgo desacoplando la selección de la acción de su evaluación:
\begin{equation}
  y_i^{\text{DDQN}} = r + \gamma (1 - d) Q\left( s', \, \arg\max_{a'} Q(s', a'; \boldsymbol{\theta}); \, \boldsymbol{\theta}^- \right)
\end{equation}
La red online $\boldsymbol{\theta}$ selecciona la mejor acción codiciosa, pero la red objetivo $\boldsymbol{\theta}^-$ evalúa objetivamente cuánto vale dicha acción.

\section{Implementación en PyTorch: Red Neuronal y Entrenamiento}
\begin{lstlisting}[style=rlcode]
import torch
import torch.nn as nn
import torch.optim as optim
import random
from collections import deque

class QNetwork(nn.Module):
    def __init__(self, state_dim=4, action_dim=4):
        super().__init__()
        self.net = nn.Sequential(
            nn.Linear(state_dim, 64),
            nn.ReLU(),
            nn.Linear(64, 64),
            nn.ReLU(),
            nn.Linear(64, action_dim)
        )
    def forward(self, x):
        return self.net(x)

class ReplayBuffer:
    def __init__(self, capacity=50000):
        self.buffer = deque(maxlen=capacity)
    def push(self, s, a, r, s_next, done):
        self.buffer.append((s, a, r, s_next, done))
    def sample(self, batch_size):
        batch = random.sample(self.buffer, batch_size)
        s, a, r, s_next, d = zip(*batch)
        return (torch.FloatTensor(s), torch.LongTensor(a).unsqueeze(1),
                torch.FloatTensor(r).unsqueeze(1), torch.FloatTensor(s_next),
                torch.FloatTensor(d).unsqueeze(1))
    def __len__(self):
        return len(self.buffer)

# Paso de entrenamiento de DQN con Target Network
def train_step(online_net, target_net, optimizer, buffer, batch_size=32, gamma=0.99):
    if len(buffer) < batch_size:
        return 0.0
    states, actions, rewards, next_states, dones = buffer.sample(batch_size)
    
    with torch.no_grad():
        # Double DQN: seleccion con online_net, evaluacion con target_net
        best_actions = online_net(next_states).argmax(1, keepdim=True)
        next_q = target_net(next_states).gather(1, best_actions)
        targets = rewards + gamma * (1 - dones) * next_q
        
    current_q = online_net(states).gather(1, actions)
    loss = nn.MSELoss()(current_q, targets)
    
    optimizer.zero_grad()
    loss.backward()
    optimizer.step()
    return loss.item()
\end{lstlisting}

\begin{notabox}
\textbf{La culminación técnica:} OptiBot dispone ahora de un cerebro artificial con aprendizaje por refuerzo profundo capaz de procesar telemetría continua y generalizar ante situaciones nunca vistas. En el capítulo final, pondremos a prueba este sistema en un proyecto integrador completo y analizaremos las fronteras que extienden estos cimientos al estado del arte contemporáneo.
\end{notabox}

\chapter{Proyecto Final OptiBot y Fronteras del Aprendizaje por Refuerzo}
\label{cap:proyecto_fronteras}

\section{El Proyecto Integrador: OptiBot en el Almacén Dinámico}
A lo largo de los capítulos precedentes, hemos acompañado a OptiBot desde la estimación de variables aleatorias elementales hasta el despliegue de redes neuronales profundas. En este proyecto final integrador, sometemos a OptiBot al desafío definitivo: la navegación en un \textbf{almacén estocástico con dinámicas no estacionarias}.

\begin{optibot}
El entorno de simulación incorpora las siguientes complejidades del mundo real:
\begin{itemize}
  \item \textbf{Obstáculos dinámicos:} Montacargas tripulados por humanos que cruzan periódicamente los pasillos principales, bloqueando la visibilidad del LiDAR.
  \item \textbf{Fricción irregular:} Derrames ocasionales de lubricante industrial que incrementan la probabilidad de que una orden de avance resulte en un deslizamiento lateral.
  \item \textbf{Picos de demanda:} Pedidos urgentes que otorgan bonificaciones de recompensa temporales si se entregan antes de un número límite de pasos.
  \item \textbf{Degradación de batería no lineal:} Consumo acelerado de energía cuando el robot transporta cargas pesadas o acelera bruscamente.
\end{itemize}
\end{optibot}

\section{Benchmark Comparativo de Paradigmas}
Para evaluar objetivamente el progreso metodológico alcanzado en el curso, sometemos a prueba tres arquitecturas diferentes bajo las mismas condiciones exactas durante 1,000 episodios de evaluación:

\begin{table}[htbp]
  \centering
  \small
  \begin{tabular}{@{}p{3.8cm}p{2.8cm}p{3.2cm}p{4.2cm}@{}}
    \toprule
    \textbf{Agente} & \textbf{Retorno Medio} & \textbf{Tasa de Éxito} & \textbf{Comportamiento Observado} \\
    \midrule
    \textbf{1. Política Aleatoria} & $-45.2 \pm 8.1$ & $2.4\%$ & Movimiento errático browniano; agota batería sin encontrar la meta. \\
    \addlinespace
    \textbf{2. Q-Learning Tabular} & $+28.6 \pm 12.3$ & $71.8\%$ & Aprende la ruta directa ideal en condiciones fijas, pero fracasa si un montacargas bloquea su pasillo memorizado. \\
    \addlinespace
    \textbf{3. Deep Q-Network (DQN)} & $\mathbf{+47.8 \pm 4.2}$ & $\mathbf{94.6\%}$ & Generaliza a partir de las distancias continuas del LiDAR; rodea obstáculos móviles y preserva batería. \\
    \bottomrule
  \end{tabular}
  \caption{Resultados del benchmark comparativo en el Almacén Inteligente de OptiBot.}
  \label{tab:benchmark}
\end{table}

\section{Diagnóstico y Depuración de Algoritmos de RL}
El entrenamiento de agentes de aprendizaje por refuerzo es notoriamente más difícil de depurar que el aprendizaje supervisado convencional: si la pérdida de una red neuronal supervisada no disminuye, sabemos que el modelo o la tasa de aprendizaje fallan; en RL, el agente genera sus propios datos, por lo que una política deficiente genera malos datos que a su vez impiden aprender una buena política.

A continuación presentamos la guía de diagnóstico para las patologías más habituales:
\begin{enumerate}
  \item \textbf{Colapso Prematuro de Exploración:} Si la tasa $\varepsilon$ decae demasiado rápido (por ejemplo, $\varepsilon_{t+1} = 0.90 \cdot \varepsilon_t$), el robot empezará a explotar su conocimiento cuando apenas ha recorrido el $5\%$ del almacén. Quedará atrapado en un óptimo local mediocre (como dar vueltas en círculo en una esquina segura).
  \item \textbf{Recompensas Excesivamente Dispersas (\textit{Sparse Rewards}):} Si el entorno entrega una recompensa de $0$ en cada paso y sólo $+100$ cuando el robot entrega el paquete, un agente aleatorio tardará millones de pasos en toparse con la meta por casualidad. La solución consiste en aplicar \textbf{Reward Shaping} (por ejemplo, conceder una pequeña recompensa proporcional a la reducción de distancia euclidiana hacia la meta).
  \item \textbf{Reward Hacking (La Trampa de la Recompensa):} El agente maximiza la fórmula matemática exacta que programaste, no lo que tú tenías en mente. Si penalizas la colisión con $-10$ y no penalizas el tiempo transcurrido, OptiBot aprenderá que la mejor forma de no perder puntos es \textbf{quedarse completamente inmóvil}.
\end{enumerate}

\section{Las Fronteras del Aprendizaje por Refuerzo Moderno}
El dominio de los fundamentos cubiertos en este texto abre las puertas directas a las áreas de investigación más activas de la inteligencia artificial contemporánea:

\subsection{1. Métodos de Gradiente de Política (\textit{Policy Gradients})}
En problemas donde las acciones son continuas (como aplicar voltajes precisos a los motores de las ruedas o controlar los 7 grados de libertad de un brazo robótico articulado), discretizar el espacio de acciones resulta ineficiente. Algoritmos como \textbf{REINFORCE}, \textbf{Actor-Critic (A2C/A3C)} y \textbf{PPO} (\textit{Proximal Policy Optimization}) parametrizan directamente la política $\pi_\theta(a \mid s)$ y optimizan el retorno esperado mediante el \textbf{Teorema del Gradiente de Política}:
\begin{equation}
  \nabla_\theta J(\theta) = \mathbb{E}_{\pi_\theta} \left[ \nabla_\theta \log \pi_\theta(a \mid s) \, Q^{\pi_\theta}(s, a) \right]
\end{equation}

\subsection{2. Aprendizaje por Refuerzo a partir de Retroalimentación Humana (RLHF)}
La explosión de los Modelos de Lenguaje Masivos (LLMs) como ChatGPT, Claude y Gemini se sustenta en el Aprendizaje por Refuerzo. Tras entrenar un modelo base con billones de palabras para predecir el siguiente token (aprendizaje auto-supervisado), el modelo se alinea con las preferencias humanas mediante un proceso de tres etapas:
\begin{enumerate}
  \item Evaluadores humanos puntúan y comparan pares de respuestas generadas por el modelo.
  \item Se entrena una red neuronal que actúa como \textbf{Modelo de Recompensa} (\textit{Reward Model}), aprendiendo a predecir qué respuestas agradan más al usuario.
  \item Se utiliza el algoritmo de RL \textbf{PPO} para ajustar los pesos del LLM (el agente), donde la acción es emitir la siguiente palabra y la recompensa es la calificación asignada por el Modelo de Recompensa.
\end{enumerate}

\section{Epílogo del Curso}
Has recorrido el camino completo: desde la definición de un espacio muestral y el cálculo manual de esperanzas matemáticas, pasando por las Cadenas de Markov, la ecuación lineal de Bellman, la Programación Dinámica, los métodos Monte Carlo, el error de Diferencia Temporal en SARSA y Q-Learning, hasta las redes neuronales profundas de DQN y los fundamentos de la alineación de LLMs con RLHF.

Cada ecuación avanzada del estado del arte contemporáneo no es más que la refinación lógica de los mismos tres pilares que ahora dominas con total claridad: \textbf{Probabilidad, Ecuaciones de Bellman y Aprendizaje por Bootstrapping}.


\end{document}
